\documentclass[11pt,oneside,openany]{book}

% ---------- Encoding and typography ----------
\usepackage[T1]{fontenc}
\usepackage[utf8]{inputenc}
\usepackage{lmodern}
\usepackage{microtype}
\usepackage[letterpaper,margin=0.82in,headheight=15pt]{geometry}
\usepackage{setspace}
\setstretch{1.06}
\usepackage{ragged2e}

% ---------- Mathematics ----------
\usepackage{amsmath,amssymb,amsthm,mathtools,bm}
\usepackage{dsfont}
\usepackage{cancel}

% ---------- Tables and figures ----------
\usepackage{booktabs,array,tabularx,longtable,multirow,makecell}
\usepackage{graphicx}
\usepackage{float}
\usepackage{caption}
\usepackage{subcaption}
\usepackage{tikz}
\usetikzlibrary{arrows.meta,positioning,calc,fit,backgrounds,matrix,shapes.geometric,decorations.pathreplacing,patterns,3d}
\usepackage{pgfplots}
\pgfplotsset{compat=1.18}

% ---------- Boxes and code ----------
\usepackage[most]{tcolorbox}
\usepackage{listings}
\usepackage{enumitem}

% ---------- Navigation and section styling ----------
\usepackage{fancyhdr}
\usepackage{titlesec}
\usepackage{hyperref}
\usepackage[nameinlink,noabbrev]{cleveref}
\usepackage{bookmark}

% ---------- Color system ----------
\definecolor{navy}{HTML}{15324A}
\definecolor{teal}{HTML}{167D7F}
\definecolor{sky}{HTML}{DCEFF2}
\definecolor{gold}{HTML}{D99A2B}
\definecolor{amber}{HTML}{FFF3D6}
\definecolor{brick}{HTML}{A83F39}
\definecolor{rose}{HTML}{F7E4E2}
\definecolor{ink}{HTML}{20262D}
\definecolor{slate}{HTML}{55616D}
\definecolor{mist}{HTML}{F4F7F9}
\definecolor{green}{HTML}{2C7A4B}
\definecolor{lightgreen}{HTML}{E7F3EB}
\definecolor{purple}{HTML}{7257A6}
\definecolor{lightpurple}{HTML}{EEE9F8}

\color{ink}

% ---------- PDF metadata ----------
\hypersetup{
  colorlinks=true,
  linkcolor=navy,
  citecolor=teal,
  urlcolor=brick,
  pdftitle={Contextual Similarity Optimization for Robust Pose-Sequence Retrieval},
  pdfauthor={Prepared as a study guide for Giacomo Cappelletto},
  pdfsubject={Mathematics and implementation of contextual metric learning},
  pdfkeywords={metric learning, contextual similarity, pose retrieval, MotionBERT, nearest neighbors}
}

% ---------- Headers and chapter styling ----------
\pagestyle{fancy}
\fancyhf{}
\fancyhead[L]{\small\textcolor{slate}{Contextual Similarity Optimization}}
\fancyhead[R]{\small\textcolor{slate}{Chapter \thechapter}}
\fancyfoot[C]{\small\thepage}
\renewcommand{\headrulewidth}{0.3pt}

\titleformat{\chapter}[display]
  {\normalfont\bfseries\color{navy}}
  {\filleft\fontsize{42}{44}\selectfont\thechapter}
  {0.2ex}
  {\titlerule[1.2pt]\vspace{1ex}\Huge}
  [\vspace{0.5ex}\titlerule]
\titlespacing*{\chapter}{0pt}{-20pt}{24pt}
\titleformat{\section}{\Large\bfseries\color{navy}}{\thesection}{0.8em}{}
\titleformat{\subsection}{\large\bfseries\color{teal}}{\thesubsection}{0.7em}{}

% ---------- Theorem styles ----------
\newtheoremstyle{myplain}{8pt}{8pt}{\itshape}{}{}{.}{0.5em}{}
\theoremstyle{myplain}
\newtheorem{proposition}{Proposition}[chapter]
\newtheorem{lemma}[proposition]{Lemma}
\theoremstyle{definition}
\newtheorem{definition}[proposition]{Definition}
\newtheorem{example}[proposition]{Example}

% ---------- Custom boxes ----------
\newtcolorbox{conceptbox}[1][]{
  enhanced,breakable,colback=sky,colframe=teal,boxrule=0.8pt,
  arc=2mm,left=2.5mm,right=2.5mm,top=2mm,bottom=2mm,
  title=Concept,#1,fonttitle=\bfseries\color{navy}}
\newtcolorbox{paperbox}[1][]{
  enhanced,breakable,colback=mist,colframe=navy,boxrule=0.8pt,
  arc=2mm,left=2.5mm,right=2.5mm,top=2mm,bottom=2mm,
  title=Original paper,#1,fonttitle=\bfseries\color{navy}}
\newtcolorbox{implementationbox}[1][]{
  enhanced,breakable,colback=lightgreen,colframe=green,boxrule=0.8pt,
  arc=2mm,left=2.5mm,right=2.5mm,top=2mm,bottom=2mm,
  title=Implementation,#1,fonttitle=\bfseries\color{green}}
\newtcolorbox{warningbox}[1][]{
  enhanced,breakable,colback=rose,colframe=brick,boxrule=0.8pt,
  arc=2mm,left=2.5mm,right=2.5mm,top=2mm,bottom=2mm,
  title=Do not overclaim,#1,fonttitle=\bfseries\color{brick}}
\newtcolorbox{meetingbox}[1][]{
  enhanced,breakable,colback=amber,colframe=gold,boxrule=0.9pt,
  arc=2mm,left=2.5mm,right=2.5mm,top=2mm,bottom=2mm,
  title=For the discussion,#1,fonttitle=\bfseries\color{navy}}
\newtcolorbox{derivationbox}[1][]{
  enhanced,breakable,colback=lightpurple,colframe=purple,boxrule=0.8pt,
  arc=2mm,left=2.5mm,right=2.5mm,top=2mm,bottom=2mm,
  title=Derivation,#1,fonttitle=\bfseries\color{purple}}

% ---------- Listings ----------
\lstdefinestyle{pythonstyle}{
  language=Python,
  basicstyle=\ttfamily\footnotesize,
  keywordstyle=\color{purple}\bfseries,
  commentstyle=\color{green},
  stringstyle=\color{brick},
  numbers=left,
  numberstyle=\tiny\color{slate},
  numbersep=7pt,
  frame=single,
  rulecolor=\color{slate!50},
  backgroundcolor=\color{mist},
  showstringspaces=false,
  breaklines=true,
  columns=fullflexible,
  tabsize=4,
  xleftmargin=1em,
  framexleftmargin=0.7em
}
\lstset{style=pythonstyle}

% ---------- Mathematical notation ----------
\newcommand{\R}{\mathbb{R}}
\newcommand{\E}{\mathbb{E}}
\newcommand{\ind}{\mathds{1}}
\newcommand{\sg}{\operatorname{sg}}
\newcommand{\simcos}{\operatorname{cos}}
\newcommand{\relu}{[\,\cdot\,]_+}
\newcommand{\norm}[1]{\left\lVert #1\right\rVert}
\newcommand{\set}[1]{\left\{#1\right\}}
\newcommand{\card}[1]{\left|#1\right|}
\newcommand{\inner}[2]{\left\langle #1,#2\right\rangle}
\newcommand{\transpose}{\mathsf{T}}
\newcommand{\had}{\odot}
\newcommand{\paperEq}[1]{\tag{P#1}}
\newcommand{\code}[1]{\texttt{#1}}
\newcommand{\Y}{\mathbf{Y}}
\newcommand{\Smat}{\mathbf{S}}
\newcommand{\Dmat}{\mathbf{D}}
\newcommand{\Nmat}{\mathbf{N}}
\newcommand{\Wmat}{\mathbf{W}}

% ---------- TikZ styles ----------
\tikzset{
  block/.style={rounded corners=2mm,draw=navy,thick,fill=mist,minimum height=9mm,align=center,inner sep=4pt},
  data/.style={rounded corners=2mm,draw=teal,thick,fill=sky,minimum height=9mm,align=center,inner sep=4pt},
  loss/.style={rounded corners=2mm,draw=brick,thick,fill=rose,minimum height=9mm,align=center,inner sep=4pt},
  impl/.style={rounded corners=2mm,draw=green,thick,fill=lightgreen,minimum height=9mm,align=center,inner sep=4pt},
  arrow/.style={-{Latex[length=2.5mm]},thick,draw=slate},
  dashedarrow/.style={-{Latex[length=2.5mm]},thick,dashed,draw=slate}
}

\begin{document}
\hypersetup{pageanchor=false}

% ================================================================
% Cover
% ================================================================
\begin{titlepage}
\begin{tikzpicture}[remember picture,overlay]
  \fill[navy] (current page.north west) rectangle ([yshift=-3.65in]current page.north east);
  \fill[teal] ([yshift=-3.65in]current page.north west) rectangle ([yshift=-3.83in]current page.north east);
  \draw[gold,line width=3pt] ([xshift=0.85in,yshift=-0.88in]current page.north west) -- ([xshift=-0.85in,yshift=-0.88in]current page.north east);
\end{tikzpicture}

\vspace*{0.55in}
{\color{white}\sffamily\bfseries\fontsize{28}{33}\selectfont
Contextual Similarity Optimization\\[0.14in]
for Robust Pose-Sequence Retrieval\par}
\vspace{0.32in}
{\color{white}\sffamily\Large
A beginner-first mathematical and implementation guide to\\
Liao, Tsiligkaridis, and Kulis (ICML 2023)\par}

\vspace{1.42in}
\begin{center}
\resizebox{0.97\textwidth}{!}{%
\begin{tikzpicture}[node distance=9mm]
  \node[data,minimum width=29mm] (x) {pose sequence\\$x\in\R^{T\times J\times C}$};
  \node[block,right=of x,minimum width=29mm] (enc) {temporal\\encoder};
  \node[data,right=of enc,minimum width=25mm] (z) {unit embedding\\$z\in\mathbb S^{d-1}$};
  \node[block,right=of z,minimum width=27mm] (ctx) {contextual\\neighborhoods};
  \node[loss,right=of ctx,minimum width=27mm] (ret) {robust\\retrieval};
  \draw[arrow] (x)--(enc);
  \draw[arrow] (enc)--(z);
  \draw[arrow] (z)--(ctx);
  \draw[arrow] (ctx)--(ret);
\end{tikzpicture}%
}
\end{center}

\vfill
\begin{minipage}{0.9\textwidth}
\small
\textbf{Audience.} Begins below the level of the paper: variables, indices, indicators, vectors, matrices, norms, dot products, and losses are introduced explicitly. It then develops the paper's equations from first principles, derives the neighborhood objective, explains the non-standard gradients, and maps the method into a controlled pose-retrieval experiment.
\end{minipage}

\vspace{0.28in}
{\large Expanded beginner-first study edition}\par
{\small July 2026}
\end{titlepage}
\hypersetup{pageanchor=true}

\frontmatter
\chapter*{How to use this guide today}
\addcontentsline{toc}{chapter}{How to use this guide today}

\begin{meetingbox}[title=Two routes through the book]
\textbf{Foundation route.} Read \cref{chap:math-primer,chap:pose-to-vector,chap:geometry,chap:training-primer} in order. This route begins with the meaning of indices, indicators, vectors, matrices, and losses before using the paper's notation.

\medskip
\textbf{Rapid discussion route.} Read the research claim below, then \cref{fig:full-pipeline}, the five central equations \cref{eq:cosdist,eq:neighborhood,eq:w1,eq:qe,eq:full-objective}, the ranking proposition in \cref{sec:proposition}, the implementation plan in \cref{chap:implementation}, and the discussion script in \cref{chap:meeting}.
\end{meetingbox}

\begin{conceptbox}[title=The one-sentence research claim]
\textbf{Test whether a neighborhood-aware metric-learning objective preserves the ranking quality of pose-sequence embeddings better than standard objectives when query skeletons are jittered, incomplete, temporally damaged, or produced by a shifted pose estimator.}
\end{conceptbox}

\begin{warningbox}[title=The scientific boundary]
The contextual paper demonstrates robustness to \emph{label noise and overfitting in image retrieval}. It does \textbf{not} establish robustness to corrupted pose coordinates. The proposed project is compelling precisely because that transfer is an open, falsifiable hypothesis rather than a guaranteed consequence.
\end{warningbox}

\begin{paperbox}[title=Source-fidelity note]
Equations marked \textbf{P1--P31} reproduce or algebraically restate the numbered equations in Liao, Tsiligkaridis, and Kulis. Motion-specific equations, implementation recommendations, and robustness metrics are explicitly labeled as adaptations or proposed extensions. One line in the paper's printed Algorithm 1 appears inconsistent with its Eq.~(5): the algorithm line shows an MSE between contextual and cosine similarity, while Eq.~(5), the proof, and the released code optimize contextual similarity against the binary label-similarity matrix. This guide follows Eq.~(5) and the reference implementation.
\end{paperbox}

\begin{conceptbox}[title=How every difficult equation should be studied]
For each major equation, first identify the output object, then read every symbol aloud, then compute a tiny numerical example, and only then study the formal or implementation-level interpretation. The book intentionally repeats core ideas at increasing levels of maturity.
\end{conceptbox}

\tableofcontents
\listoffigures

\mainmatter

% ================================================================
\chapter{How to Read the Mathematics in This Guide}
\label{chap:math-primer}
% ================================================================

\section{The explanation ladder used throughout the book}

This edition deliberately explains each important idea in several passes. The passes are not repetitions. Each one adds a new level of precision.

\begin{enumerate}[leftmargin=*,itemsep=4pt]
  \item \textbf{Start with the task.} What is the computation trying to accomplish?
  \item \textbf{Use a tiny concrete example.} Work with two or four samples rather than an abstract dataset.
  \item \textbf{Read the formula aloud.} Translate every symbol into a sentence.
  \item \textbf{Compute it by hand.} Substitute actual numbers.
  \item \textbf{State the formal interpretation.} Explain why the formula has the required mathematical properties.
  \item \textbf{Connect it to code.} Identify tensor shapes and operations.
\end{enumerate}

\begin{conceptbox}[title=The rule for reading a formula]
Do not begin by trying to understand the entire expression at once. First identify:
\[
\boxed{\text{what is being defined}}\quad
\boxed{\text{what information enters it}}\quad
\boxed{\text{what type of object comes out}}.
\]
Only after those three questions are answered should you inspect the individual operators.
\end{conceptbox}

\section{Scalars, vectors, matrices, sets, and functions}

Most confusion in machine-learning mathematics comes from not knowing what kind of object a symbol represents. The same-looking letter can refer to a single number, a list of numbers, or a two-dimensional table. These objects support different operations.

\begin{center}
\begin{tabularx}{\textwidth}{>{\bfseries}l l X}
\toprule
Object & Example & Meaning in this project \\
\midrule
Scalar & $s_{ij}\in\R$ & one similarity value between two samples \\
Vector & $z_i\in\R^d$ & one embedding containing $d$ numbers \\
Matrix & $S\in\R^{n\times n}$ & all pairwise similarities in a batch \\
Set & $N_k(i)$ & the collection of the $k$ nearest samples to $i$ \\
Function & $f_\theta(x)$ & a rule, implemented by a neural network, that maps an input to an embedding \\
Binary value & $y_{ij}\in\{0,1\}$ & an answer to a yes/no question \\
\bottomrule
\end{tabularx}
\end{center}

\subsection{A scalar is one number}

A scalar has no list structure. Examples are
\[
  3.2,\qquad -1,\qquad 0.87.
\]
In the contextual paper, $s_{ij}$ is a scalar. It is one number telling us how similar sample $i$ is to sample $j$.

\subsection{A vector is an ordered list of numbers}

A vector with three coordinates can be written as
\[
  z_i=\begin{bmatrix}0.6\\0.8\\0\end{bmatrix}\in\R^3.
\]
The superscript $3$ in $\R^3$ means that the vector has three real-valued coordinates. In the actual project, the dimension may be $128$, $256$, or another chosen value.

\begin{conceptbox}[title=Why order matters]
The vectors $[1,0]^{\mathsf T}$ and $[0,1]^{\mathsf T}$ contain the same two numbers, but they are not the same vector. Coordinate 1 and coordinate 2 represent different learned directions in the embedding space.
\end{conceptbox}

\subsection{A matrix is a table of numbers}

A matrix has rows and columns. For three samples, a similarity matrix may look like
\[
S=
\begin{bmatrix}
1 & 0.8 & -0.2\\
0.8 & 1 & 0.1\\
-0.2 & 0.1 & 1
\end{bmatrix}.
\]
The entry in row $i$ and column $j$ is written $S_{ij}$ or $s_{ij}$. The diagonal entries are all $1$ because every normalized vector is perfectly similar to itself.

\subsection{A set is a collection of distinct objects}

The notation
\[
  N_3(i)=\{i,2,7\}
\]
means that the three nearest items to sample $i$ are sample $i$ itself, sample $2$, and sample $7$. Curly braces describe membership, not an ordered vector. In the paper, sets are often converted to binary indicator vectors so that set operations can be implemented with matrix multiplication.

\subsection{A function is a rule}

The expression
\[
  z_i=f_\theta(x_i)
\]
means: apply the function $f_\theta$ to input $x_i$ and call the result $z_i$. The subscript $\theta$ represents all trainable parameters of the neural network. During gradient descent, those parameters change.

\section{Indices: what do $i$, $j$, and $p$ mean?}

An index is an integer used to identify an item. If a mini-batch contains $n$ samples, then
\[
  i,j,p\in\{1,2,\ldots,n\}.
\]
The symbols $i$ and $j$ are not values of the input. They are addresses. They tell us which samples are being discussed.

\begin{center}
\begin{tabular}{ccl}
\toprule
Index & Sample & Label \\
\midrule
$1$ & first pose sequence & walk \\
$2$ & second pose sequence & walk \\
$3$ & third pose sequence & jump \\
$4$ & fourth pose sequence & throw \\
\bottomrule
\end{tabular}
\end{center}

Thus $y_2$ means the class label of sample $2$, while $y_{24}$ means a relation between samples $2$ and $4$. The extra index changes the meaning of the symbol.

\section{A complete explanation of the label-similarity formula}
\label{sec:indicator-primer}

The formula that often appears first in the paper is
\begin{equation}
  y_{ij}=\ind[y_i=y_j]\in\{0,1\}.
  \label{eq:indicator-primer}
\end{equation}

We now unpack every part.

\subsection{Pass 1: the simplest possible meaning}

The formula asks one yes/no question:

\begin{center}
\Large\bfseries
Do sample $i$ and sample $j$ have the same class label?
\end{center}

If the answer is yes, assign the number $1$. If the answer is no, assign the number $0$.

\begin{figure}[H]
\centering
\begin{tikzpicture}[node distance=12mm]
  \node[data,minimum width=31mm] (a) {sample $i$\\label: walk};
  \node[data,right=of a,minimum width=31mm] (b) {sample $j$\\label: walk};
  \node[block,right=of b,minimum width=34mm] (q) {Are the labels equal?\\$y_i=y_j$};
  \node[impl,right=of q,minimum width=24mm] (o) {$y_{ij}=1$};
  \draw[arrow] (a)--(q);
  \draw[arrow] (b)--(q);
  \draw[arrow] (q)--node[above]{yes}(o);
\end{tikzpicture}
\caption{The label-similarity value is a numerical encoding of a yes/no statement.}
\end{figure}

\subsection{Pass 2: read the expression aloud}

Read
\[
  y_{ij}=\ind[y_i=y_j]\in\{0,1\}
\]
as follows:

\begin{quote}
``The pairwise label-similarity value for samples $i$ and $j$ equals the indicator of the statement that the label of sample $i$ equals the label of sample $j$. The result must be either zero or one.''
\end{quote}

\subsection{Pass 3: symbol-by-symbol dictionary}

\begin{center}
\begin{tabularx}{\textwidth}{>{\bfseries}l X}
\toprule
Symbol & Meaning \\
\midrule
$i$ & index of the first sample \\
$j$ & index of the second sample \\
$y_i$ & class label attached to sample $i$ \\
$y_j$ & class label attached to sample $j$ \\
$y_i=y_j$ & a logical statement that is either true or false \\
$\ind[\,\cdot\,]$ & indicator function: convert true to $1$ and false to $0$ \\
$y_{ij}$ & the resulting pairwise binary similarity target \\
$\in$ & ``is an element of'' or ``belongs to'' \\
$\{0,1\}$ & the set containing exactly the values zero and one \\
\bottomrule
\end{tabularx}
\end{center}

\subsection{Pass 4: the indicator function by itself}

For any logical statement $A$, define
\[
\ind[A]=
\begin{cases}
1,&\text{if $A$ is true},\\
0,&\text{if $A$ is false}.
\end{cases}
\]
For example,
\[
\ind[3<5]=1,\qquad
\ind[8=2]=0.
\]
The indicator function lets us place logical information inside arithmetic expressions. Once a statement has been converted to $0$ or $1$, we can add it, multiply it, average it, or put it into a matrix.

\subsection{Pass 5: a four-sample example}

Suppose the label vector is
\[
  \bm y=
  \begin{bmatrix}
  \text{walk}\\
  \text{walk}\\
  \text{jump}\\
  \text{throw}
  \end{bmatrix}.
\]
Then
\[
  y_{12}=\ind[\text{walk}=\text{walk}]=1,
\]
while
\[
  y_{13}=\ind[\text{walk}=\text{jump}]=0.
\]
Applying the rule to every ordered pair gives
\[
Y=
\begin{bmatrix}
1&1&0&0\\
1&1&0&0\\
0&0&1&0\\
0&0&0&1
\end{bmatrix}.
\]

\begin{conceptbox}[title=Why the matrix is symmetric]
If sample $i$ has the same label as sample $j$, then sample $j$ also has the same label as sample $i$. Therefore
\[
y_{ij}=y_{ji},
\]
so $Y=Y^{\mathsf T}$.
\end{conceptbox}

\begin{conceptbox}[title=Why the diagonal is one]
Every label equals itself. Therefore $y_i=y_i$ is always true and
\[
y_{ii}=1.
\]
The contextual loss later removes diagonal terms because self-similarity is trivial and does not teach the model how to rank different samples.
\end{conceptbox}

\subsection{Pass 6: the mature interpretation}

The label-similarity matrix $Y$ is the adjacency matrix of a graph. Each sample is a node. Two nodes are connected when they share a class label. Training attempts to make the learned similarity structure resemble this target graph. Contextual loss differs from a purely pairwise loss because it compares not only individual edges but also the local graph surrounding each node.

\section{Common operators used in the paper}

\subsection{Equality, definition, and membership}

The symbol $=$ may report an equality or define a new object. Context usually tells which role it has. The symbol $\in$ reports membership:
\[
  z_i\in\R^d
\]
means that $z_i$ is a $d$-dimensional real vector.

\subsection{Summation notation}

The expression
\[
  \sum_{p=1}^{n} a_p
\]
means
\[
  a_1+a_2+\cdots+a_n.
\]
The letter $p$ is a temporary index. It takes every integer value from $1$ through $n$.

For example,
\[
  \sum_{p=1}^{4} a_p
  =a_1+a_2+a_3+a_4.
\]

\subsection{Conditions underneath a summation}

A sum such as
\[
  \sum_{i,j\,|\,i\ne j}(y_{ij}-w_{ij})^2
\]
means: include all pairs $(i,j)$ except pairs for which $i=j$. The vertical bar can be read as ``such that.''

\subsection{Transpose}

For a column vector
\[
z=\begin{bmatrix}z_1\\z_2\\z_3\end{bmatrix},
\qquad
z^{\mathsf T}=\begin{bmatrix}z_1&z_2&z_3\end{bmatrix}.
\]
The transpose turns rows into columns and columns into rows. It enables matrix expressions such as $ZZ^{\mathsf T}$.

\subsection{Norm}

The Euclidean or $L_2$ norm is the length of a vector:
\[
  \norm{z}_2=\sqrt{z_1^2+z_2^2+\cdots+z_d^2}.
\]
For $z=[3,4]^{\mathsf T}$,
\[
  \norm{z}_2=\sqrt{3^2+4^2}=5.
\]

\subsection{Dot product}

For two vectors $a,b\in\R^d$,
\[
  a^{\mathsf T}b=\sum_{r=1}^{d}a_rb_r.
\]
This means: multiply matching coordinates, then add the products. For
\[
  a=\begin{bmatrix}1\\2\\3\end{bmatrix},\qquad
  b=\begin{bmatrix}4\\0\\-1\end{bmatrix},
\]
we get
\[
  a^{\mathsf T}b=(1)(4)+(2)(0)+(3)(-1)=1.
\]

\subsection{The positive-part operator}

The notation
\[
  [u]_+=\max(u,0)
\]
means that negative values are replaced by zero. It is the scalar ReLU operation. For example,
\[
  [-2]_+=0,\qquad [0.7]_+=0.7.
\]
Metric-learning margin losses use this operation so that already-correct pairs contribute no loss.

\subsection{Set intersection and complement}

If
\[
A=\{1,2,4\},\qquad B=\{2,4,5\},
\]
then their intersection is
\[
  A\cap B=\{2,4\}.
\]
The intersection contains elements that appear in both sets. Contextual similarity relies heavily on the size of neighborhood intersections.

\section{How a matrix formula becomes code}

Suppose $Z\in\R^{n\times d}$ contains one embedding per row. Then
\[
  S=ZZ^{\mathsf T}
\]
produces an $n\times n$ matrix. Entry $(i,j)$ is
\[
  S_{ij}=z_i^{\mathsf T}z_j.
\]
In PyTorch, this is simply
\begin{lstlisting}
S = Z @ Z.T
\end{lstlisting}
The symbol \code{@} means matrix multiplication. This single operation computes all $n^2$ pairwise dot products in parallel.

\begin{meetingbox}[title=Checkpoint before continuing]
You are ready for the next chapter when you can explain the difference between $y_i$ and $y_{ij}$, state what an indicator function does, and identify whether $z_i$, $S$, and $N_k(i)$ are respectively a vector, matrix, or set.
\end{meetingbox}

% ================================================================
\chapter{From a Pose Sequence to a Point in Embedding Space}
\label{chap:pose-to-vector}
% ================================================================

\section{Start with the physical data}

A pose sequence is a series of estimated body configurations. At each time step, a pose estimator returns the coordinates of several body joints. A convenient mathematical representation is
\[
  x_i\in\R^{T\times J\times C},
\]
where:
\begin{itemize}[leftmargin=*]
  \item $i$ identifies the motion clip;
  \item $T$ is the number of time frames;
  \item $J$ is the number of body joints;
  \item $C$ is the number of coordinates per joint, usually $2$ or $3$.
\end{itemize}

If a clip contains $T=120$ frames, $J=25$ joints, and three coordinates $(x,y,z)$ per joint, then the clip contains
\[
120\times25\times3=9000
\]
coordinate values before any extra confidence values or features are added.

\begin{figure}[H]
\centering
\begin{tikzpicture}[node distance=8mm]
  \node[data,minimum width=30mm,minimum height=18mm] (frames) {frames\\$t=1,\ldots,T$};
  \node[data,right=of frames,minimum width=30mm,minimum height=18mm] (joints) {joints\\$j=1,\ldots,J$};
  \node[data,right=of joints,minimum width=30mm,minimum height=18mm] (coords) {coordinates\\$c\in\{x,y,z\}$};
  \node[block,right=12mm of coords,minimum width=34mm] (tensor) {pose tensor\\$x_i\in\R^{T\times J\times C}$};
  \draw[arrow] (frames)--(tensor);
  \draw[arrow] (joints)--(tensor);
  \draw[arrow] (coords)--(tensor);
\end{tikzpicture}
\caption{A pose sequence is a three-axis array: time, joint identity, and coordinate.}
\end{figure}

\section{Why the raw tensor is not directly used for retrieval}

Two recordings of the same action can differ in length, body proportions, camera viewpoint, execution speed, and pose-estimation error. Directly subtracting all raw coordinates would therefore measure many nuisance differences. We instead learn an encoder
\[
  h_i=f_\theta(x_i)\in\R^d.
\]
The encoder compresses the entire sequence into a fixed-length vector. The training objective determines what information that vector should preserve.

\begin{figure}[H]
\centering
\resizebox{\textwidth}{!}{%
\begin{tikzpicture}[node distance=8mm]
  \node[data,minimum width=32mm] (x) {pose sequence\\$x_i$};
  \node[block,right=of x,minimum width=36mm] (enc) {MotionBERT or another\\temporal encoder $f_\theta$};
  \node[data,right=of enc,minimum width=34mm] (h) {raw embedding\\$h_i\in\R^d$};
  \node[block,right=of h,minimum width=28mm] (norm) {$L_2$\\normalization};
  \node[data,right=of norm,minimum width=34mm] (z) {unit embedding\\$z_i\in\mathbb S^{d-1}$};
  \draw[arrow] (x)--(enc);
  \draw[arrow] (enc)--(h);
  \draw[arrow] (h)--(norm);
  \draw[arrow] (norm)--(z);
\end{tikzpicture}%
}
\caption{The retrieval model converts a variable, structured motion sequence into one fixed-length point.}
\end{figure}

\section{What an embedding coordinate means}

A learned coordinate usually does not have a direct human-readable meaning. Coordinate 17 is not necessarily ``elbow angle'' and coordinate 44 is not necessarily ``speed.'' The network distributes useful information across many coordinates. What matters is the geometry of the complete vector: samples that should retrieve one another must point in similar directions.

\section{Normalization with actual numbers}

Suppose an encoder outputs
\[
  h_i=\begin{bmatrix}3\\4\end{bmatrix}.
\]
Its length is
\[
  \norm{h_i}_2=\sqrt{3^2+4^2}=5.
\]
The normalized embedding is
\[
  z_i=\frac{h_i}{\norm{h_i}_2}
  =\begin{bmatrix}3/5\\4/5\end{bmatrix}
  =\begin{bmatrix}0.6\\0.8\end{bmatrix}.
\]
Checking the new length,
\[
\norm{z_i}_2=\sqrt{0.6^2+0.8^2}=1.
\]

\begin{conceptbox}[title=What normalization changes and preserves]
Normalization changes vector magnitude but preserves vector direction. After normalization, the network cannot make two samples appear similar merely by increasing both vector lengths. Similarity depends on angle.
\end{conceptbox}

\section{The mature representation-learning view}

The model learns a map
\[
  f_\theta:\mathcal X\rightarrow\mathbb S^{d-1},
\]
from the space $\mathcal X$ of pose sequences to the unit sphere $\mathbb S^{d-1}$. Metric learning does not prescribe the coordinates of each point individually. It constrains relationships among points. The contextual paper strengthens this relational view by constraining local neighborhoods rather than only isolated pairs.


% ================================================================
\chapter{How a Loss Function Trains a Neural Network}
\label{chap:training-primer}
% ================================================================

\section{The training loop in its simplest form}

A neural network begins with parameters $\theta$ that do not yet produce useful retrieval geometry. Training repeats four operations:

\begin{enumerate}[leftmargin=*,itemsep=5pt]
  \item compute embeddings using the current parameters;
  \item measure how wrong the embeddings are using a loss function;
  \item compute how the loss changes when each parameter changes;
  \item adjust the parameters in a direction that reduces the loss.
\end{enumerate}

\begin{figure}[H]
\centering
\begin{tikzpicture}[node distance=7mm]
  \node[data,minimum width=27mm] (batch) {mini-batch\\$x_1,\ldots,x_n$};
  \node[block,right=of batch,minimum width=28mm] (model) {network\\$f_\theta$};
  \node[data,right=of model,minimum width=27mm] (emb) {embeddings\\$Z$};
  \node[loss,right=of emb,minimum width=28mm] (loss) {loss\\$\mathcal L$};
  \node[block,below=11mm of loss,minimum width=29mm] (grad) {gradients\\$\nabla_\theta\mathcal L$};
  \node[impl,left=of grad,minimum width=33mm] (update) {parameter update\\$\theta\leftarrow\theta-\eta\nabla_\theta\mathcal L$};
  \draw[arrow] (batch)--(model);
  \draw[arrow] (model)--(emb);
  \draw[arrow] (emb)--(loss);
  \draw[arrow] (loss)--(grad);
  \draw[arrow] (grad)--(update);
  \draw[arrow] (update.west) -| node[pos=0.25,below]{repeat} (model.south);
\end{tikzpicture}
\caption{Training repeatedly measures an error and changes the network parameters to reduce it.}
\end{figure}

\section{What a loss function is}

A loss function returns one scalar. Smaller is better. A basic squared-error loss is
\[
  \ell=(t-p)^2,
\]
where $t$ is a target and $p$ is a prediction.

If $t=1$ and $p=0.8$, then
\[
\ell=(1-0.8)^2=0.04.
\]
If $p=0.2$, then
\[
\ell=(1-0.2)^2=0.64.
\]
The second prediction is farther from the target and receives a larger penalty.

\section{Why use the square?}

The square has three useful effects:
\begin{itemize}[leftmargin=*]
  \item negative and positive errors both become nonnegative;
  \item large errors receive disproportionately larger penalties;
  \item the function is smooth and easy to differentiate.
\end{itemize}

\section{A derivative is a local sensitivity}

For a scalar function $\ell(p)$, the derivative
\[
  \frac{d\ell}{dp}
\]
answers: if $p$ changes by a very small amount, how quickly does the loss change?

For
\[
  \ell(p)=(t-p)^2,
\]
the derivative is
\[
  \frac{d\ell}{dp}=2(p-t).
\]
If $t=1$ and $p=0.2$, then
\[
  \frac{d\ell}{dp}=2(0.2-1)=-1.6.
\]
The negative sign says that increasing $p$ will decrease the loss.

\section{Gradient descent}

For many parameters, collect all partial derivatives into the gradient
\[
  \nabla_\theta\mathcal L.
\]
The basic update is
\[
  \theta\leftarrow\theta-\eta\nabla_\theta\mathcal L,
\]
where $\eta>0$ is the learning rate. Subtracting the gradient moves parameters locally downhill.

\begin{conceptbox}[title=The role of the loss]
The loss does not directly move embeddings. It creates gradients. Those gradients alter network parameters. On the next forward pass, the altered network produces different embeddings.
\end{conceptbox}

\section{The chain rule in the current project}

The contextual loss depends on contextual similarity, which depends on neighborhoods, which depend on distances, which depend on cosine similarities, which depend on embeddings, which depend on network parameters:
\[
\theta
\longrightarrow Z
\longrightarrow S
\longrightarrow D
\longrightarrow N
\longrightarrow W
\longrightarrow\mathcal L.
\]
The chain rule propagates sensitivity backward through this sequence. Schematically,
\[
\frac{\partial\mathcal L}{\partial\theta}
=
\frac{\partial\mathcal L}{\partial W}
\frac{\partial W}{\partial N}
\frac{\partial N}{\partial D}
\frac{\partial D}{\partial S}
\frac{\partial S}{\partial Z}
\frac{\partial Z}{\partial\theta}.
\]
This expression is not ordinary scalar multiplication in implementation; it summarizes a sequence of Jacobian-vector products performed by automatic differentiation.

\section{Why discrete operations cause a problem}

A hard nearest-neighbor membership decision changes suddenly from $0$ to $1$. Such a decision is constant almost everywhere. Its true derivative is zero almost everywhere, so ordinary backpropagation cannot tell the model how to change a distance to cross the membership boundary.

The contextual paper keeps the exact binary decision in the forward pass but defines a heuristic nonzero derivative in the backward pass. This is the key optimization device studied in \cref{chap:step1}.

\begin{meetingbox}[title=Checkpoint]
The mature point to remember is that contextual similarity is not difficult because its forward calculation is impossible. It is difficult because the calculation contains discrete neighborhood decisions that block ordinary gradients.
\end{meetingbox}


% ================================================================
\chapter{Retrieval Geometry from First Principles}
\label{chap:geometry}
% ================================================================

\section{Retrieval is a ranking problem}

\subsection{Start with the practical task}

Assume that a database contains many pose sequences. A new query sequence arrives. The system should return the database sequences in an order such that the most relevant examples appear first.

This output is different from classification. A classifier returns one label, such as ``throw.'' A retrieval system returns an ordered list of examples:
\[
\text{query}\longrightarrow
\bigl[	ext{most similar},\text{second most similar},\ldots\bigr].
\]
The ordering itself is the prediction.

\begin{figure}[H]
\centering
\begin{tikzpicture}[node distance=6mm]
  \node[data,minimum width=25mm] (q) {query\\``throw''};
  \node[block,right=12mm of q,minimum width=25mm] (e) {encoder\\$f_\theta$};
  \node[data,right=12mm of e,minimum width=22mm] (zq) {$z_q$};
  \node[block,right=12mm of zq,minimum width=26mm] (rank) {similarity\\ranking};
  \node[right=9mm of rank,align=left] (list) {
  \begin{tabular}{c@{\quad}l@{\quad}c}
  1 & throw & $\checkmark$\\
  2 & punch & $\times$\\
  3 & throw & $\checkmark$\\
  4 & kick & $\times$\\
  5 & throw & $\checkmark$
  \end{tabular}};
  \draw[arrow] (q)--(e);
  \draw[arrow] (e)--(zq);
  \draw[arrow] (zq)--(rank);
  \draw[arrow] (rank)--(list);
  \node[below=8mm of rank,text=slate,align=center] {gallery embeddings are computed in advance};
\end{tikzpicture}
\caption{A retrieval model predicts an ordering of database items.}
\label{fig:retrieval}
\end{figure}

\subsection{Formal notation for the gallery}

Let the query be $q$. Let the gallery be the set
\[
  \mathcal G=\{g_1,g_2,\ldots,g_m\}.
\]
Here $m$ is the number of gallery items. The retrieval system assigns one score to each pair $(q,g_j)$ and sorts the items by score.

An ordering can be represented by a permutation $\pi$. The ranked output is
\[
  g_{\pi(1)},g_{\pi(2)},\ldots,g_{\pi(m)}.
\]
The function $\pi$ simply records which original gallery index occupies each rank. For example, if the original items $g_4,g_2,g_1,g_3$ are returned in that order, then
\[
\pi(1)=4,\quad \pi(2)=2,\quad \pi(3)=1,\quad \pi(4)=3.
\]

\subsection{Relevant and irrelevant items}

For the first NTU120 experiment, two sequences can be declared relevant when they share the same action class. Define
\[
  \operatorname{rel}(q,g_j)=\ind[y_q=y_j].
\]
This is the same indicator construction introduced in \cref{sec:indicator-primer}.

\begin{conceptbox}[title=What this relevance definition does not claim]
Class-level relevance means that two clips belong to the same action category. It does not mean that they have identical biomechanics, timing, quality, or technique. NTU120 provides an objective first benchmark, not a final definition of fine-grained motion similarity.
\end{conceptbox}

\section{Similarity and distance are two ways to order points}

\subsection{Similarity: larger means closer}

A similarity score is designed so that larger values indicate a stronger match. If
\[
  s(q,g_1)=0.92,\qquad s(q,g_2)=0.31,
\]
then $g_1$ should rank ahead of $g_2$.

\subsection{Distance: smaller means closer}

A distance is designed so that smaller values indicate a stronger match. If
\[
  D(q,g_1)=0.16,\qquad D(q,g_2)=1.38,
\]
then again $g_1$ should rank ahead of $g_2$.

The paper uses both languages. It uses cosine similarity to describe retrieval and squared Euclidean distance to define neighborhood thresholds. Because embeddings are normalized, the two produce exactly the same ranking.

\section{Embeddings and normalization}

\subsection{The raw network output}

Let
\[
  h_i=f_\theta(x_i)\in\R^d.
\]
Read this as:

\begin{quote}
``The neural network with parameters $\theta$ processes input sample $x_i$ and returns a $d$-dimensional vector called $h_i$.''
\end{quote}

The vector $h_i$ may have any nonzero length. The paper normalizes it:
\begin{equation}
  z_i=\frac{h_i}{\norm{h_i}_2},\qquad \norm{z_i}_2=1.
  \label{eq:normalize}
\end{equation}

\subsection{Read the normalization formula symbol by symbol}

\begin{center}
\begin{tabularx}{\textwidth}{>{\bfseries}l X}
\toprule
Part & Meaning \\
\midrule
$h_i$ & raw embedding of sample $i$ \\
$\norm{h_i}_2$ & Euclidean length of that raw embedding \\
$h_i/\norm{h_i}_2$ & divide every coordinate by the same positive length \\
$z_i$ & resulting normalized embedding \\
$\norm{z_i}_2=1$ & the normalized embedding lies on the unit sphere \\
\bottomrule
\end{tabularx}
\end{center}

\subsection{Why normalization is important}

Without normalization, a dot product is affected by both direction and magnitude:
\[
  h_i^{\mathsf T}h_j=\norm{h_i}_2\norm{h_j}_2\cos\phi.
\]
The angle $\phi$ captures directional alignment, but the two lengths also multiply the score. After normalization, both lengths equal one:
\[
  z_i^{\mathsf T}z_j=\cos\phi.
\]
This makes the score easier to interpret and keeps it inside a fixed interval.

\section{The dot product and cosine similarity}

For normalized embeddings, the paper defines
\begin{equation}
  s_{ij}=\inner{z_i}{z_j}=z_i^{\mathsf T}z_j\in[-1,1].
  \label{eq:cosine}
\end{equation}

\subsection{Read the formula aloud}

\begin{quote}
``The cosine similarity between samples $i$ and $j$ equals the dot product of their normalized embedding vectors. It must lie between negative one and positive one.''
\end{quote}

\subsection{Compute a dot product by hand}

Let
\[
  z_i=\begin{bmatrix}0.6\\0.8\end{bmatrix},
  \qquad
  z_j=\begin{bmatrix}0.8\\0.6\end{bmatrix}.
\]
Then
\begin{align*}
  s_{ij}
  &=z_i^{\mathsf T}z_j\\
  &=(0.6)(0.8)+(0.8)(0.6)\\
  &=0.48+0.48\\
  &=0.96.
\end{align*}
The value is close to one, so the vectors point in similar directions.

\subsection{Interpret the range $[-1,1]$}

\begin{center}
\begin{tabularx}{0.88\textwidth}{c X}
\toprule
Value & Geometric meaning \\
\midrule
$s_{ij}=1$ & same direction; maximum cosine similarity \\
$s_{ij}=0$ & perpendicular directions; no directional alignment \\
$s_{ij}=-1$ & opposite directions; minimum cosine similarity \\
\bottomrule
\end{tabularx}
\end{center}

\begin{figure}[H]
\centering
\begin{tikzpicture}[scale=1.15]
  \draw[slate!60] (0,0) circle (2);
  \draw[->,slate] (-2.25,0)--(2.35,0);
  \draw[->,slate] (0,-2.25)--(0,2.35);
  \draw[->,very thick,teal] (0,0)--(1.75,0.97) node[above right] {$z_i$};
  \draw[->,very thick,gold] (0,0)--(1.13,1.65) node[above] {$z_j$};
  \draw[->,very thick,brick] (0,0)--(-1.75,0.97) node[above left] {$z_k$};
  \draw[navy,thick] (0.65,0.36) arc[start angle=29,end angle=56,radius=0.75];
  \node[navy] at (0.89,0.77) {$\phi$};
  \node[align=left,anchor=west] at (2.65,1.15) {$s_{ij}=\cos\phi$\\small angle $\Rightarrow$ high similarity};
  \node[align=left,anchor=west] at (2.65,-0.35) {$s_{ik}<0$\\opposite directions};
\end{tikzpicture}
\caption{Unit-normalized embeddings turn similarity into angular geometry.}
\label{fig:unit-sphere}
\end{figure}

\section{Why squared Euclidean distance becomes $2-2s_{ij}$}

This identity is used repeatedly in the paper:
\begin{equation}
  \boxed{D(i,j)=\norm{z_i-z_j}_2^2=2-2s_{ij}.}
  \label{eq:cosdist}
\end{equation}

\subsection{Pass 1: the result in plain language}

For unit-length embeddings, knowing cosine similarity is enough to know squared Euclidean distance. A larger similarity always produces a smaller distance.

\subsection{Pass 2: explain the left-hand side}

The vector difference
\[
  z_i-z_j
\]
points from one embedding to the other. Its norm is the straight-line distance between them. Squaring the norm removes the square root and is convenient for algebra.

\subsection{Pass 3: derive every algebraic step}

Start from the definition:
\begin{align}
D(i,j)
  &=\norm{z_i-z_j}_2^2. \label{eq:dist-step1}
\end{align}
A squared norm can be written as a vector dotted with itself:
\begin{align}
  \norm{z_i-z_j}_2^2
  &=(z_i-z_j)^{\mathsf T}(z_i-z_j). \label{eq:dist-step2}
\end{align}
Distribute the multiplication:
\begin{align}
(z_i-z_j)^{\mathsf T}(z_i-z_j)
  &=z_i^{\mathsf T}z_i-z_i^{\mathsf T}z_j-z_j^{\mathsf T}z_i+z_j^{\mathsf T}z_j. \label{eq:dist-step3}
\end{align}
The middle two scalar dot products are equal:
\[
  z_i^{\mathsf T}z_j=z_j^{\mathsf T}z_i=s_{ij}.
\]
The first and last terms are squared lengths:
\[
  z_i^{\mathsf T}z_i=\norm{z_i}_2^2=1,
  \qquad
  z_j^{\mathsf T}z_j=\norm{z_j}_2^2=1.
\]
Substituting gives
\begin{align*}
D(i,j)
  &=1-s_{ij}-s_{ij}+1\\
  &=2-2s_{ij}.
\end{align*}

\subsection{Pass 4: check the endpoints}

If the vectors are identical, $s_{ij}=1$:
\[
  D(i,j)=2-2(1)=0.
\]
If the vectors are opposite, $s_{ij}=-1$:
\[
  D(i,j)=2-2(-1)=4.
\]
Therefore $D(i,j)\in[0,4]$.

\subsection{Pass 5: why the rankings are identical}

The function
\[
  D=2-2s
\]
is strictly decreasing in $s$. If $s_{i1}>s_{i2}$, then
\[
  2-2s_{i1}<2-2s_{i2}.
\]
Thus sorting by decreasing cosine similarity is identical to sorting by increasing squared Euclidean distance.

\begin{conceptbox}[title=What changes and what does not]
The numerical scale changes: similarity lies in $[-1,1]$, while squared distance lies in $[0,4]$. The order of neighbors does not change.
\end{conceptbox}

\section{The label-similarity matrix}

For a mini-batch of $n$ labeled examples, define
\begin{equation}
  y_{ij}=\ind[y_i=y_j]\in\{0,1\}.
  \label{eq:labelsim}
\end{equation}
This is the same formula studied in detail in \cref{sec:indicator-primer}. It creates the target relation for every pair of samples.

\subsection{From one pair to a full matrix}

Collect all values into
\[
  Y=[y_{ij}]\in\{0,1\}^{n\times n}.
\]
Row $i$ describes which batch samples are relevant to sample $i$. Column $j$ describes which anchors regard sample $j$ as relevant.

\begin{figure}[H]
\centering
\begin{tikzpicture}[scale=0.62]
  \foreach \r in {0,...,11}{
    \foreach \c in {0,...,11}{
      \pgfmathtruncatemacro{\same}{(int(\r/4)==int(\c/4))?1:0}
      \ifnum\same=1
        \fill[teal!75] (\c,-\r) rectangle ++(1,-1);
      \else
        \fill[mist] (\c,-\r) rectangle ++(1,-1);
      \fi
      \draw[white,line width=0.25pt] (\c,-\r) rectangle ++(1,-1);
    }
  }
  \draw[navy,thick] (0,0) rectangle (12,-12);
  \node[left] at (0,-2) {class A};
  \node[left] at (0,-6) {class B};
  \node[left] at (0,-10) {class C};
  \node[above] at (2,0) {A};
  \node[above] at (6,0) {B};
  \node[above] at (10,0) {C};
  \node[right=8mm,align=left] at (12,-6) {teal: $y_{ij}=1$\\white: $y_{ij}=0$};
\end{tikzpicture}
\caption{With samples grouped by class, the target similarity matrix has visible blocks.}
\label{fig:label-matrix}
\end{figure}

\subsection{The learned matrix and the target matrix}

The learned cosine similarities form
\[
  S=[s_{ij}]\in[-1,1]^{n\times n}.
\]
The target relations form
\[
  Y=[y_{ij}]\in\{0,1\}^{n\times n}.
\]
A pairwise loss compares entries of $S$ with pair labels. The contextual loss first transforms $S$ into a new matrix $W$ that describes neighborhood agreement, and then compares $W$ with $Y$.

\begin{figure}[H]
\centering
\begin{tikzpicture}[node distance=12mm]
  \node[data,minimum width=28mm] (S) {learned pairwise\\similarity $S$};
  \node[block,right=of S,minimum width=34mm] (T) {construct local\\neighborhood graph};
  \node[data,right=of T,minimum width=31mm] (W) {contextual\\similarity $W$};
  \node[loss,right=of W,minimum width=28mm] (Y) {compare with\\target graph $Y$};
  \draw[arrow] (S)--(T);
  \draw[arrow] (T)--(W);
  \draw[arrow] (W)--(Y);
\end{tikzpicture}
\caption{The main contribution is the intermediate transformation from pairwise geometry to contextual geometry.}
\end{figure}

\begin{meetingbox}[title=What to say in the discussion]
``Each sequence becomes a unit vector. The dot product is its cosine similarity. Because the vectors have unit length, squared Euclidean distance is exactly $2-2s_{ij}$, so similarity and distance produce the same neighbor ranking. The labels create a binary target matrix $Y$. The contextual method does not match $S$ to $Y$ directly; it first asks whether pairs occupy similar local neighborhoods.''
\end{meetingbox}


% ================================================================
\chapter{Metric-Learning Baselines and Balanced Batches}
\label{chap:baselines}
% ================================================================

\section{Begin with the purpose of a mini-batch}

A dataset may contain tens of thousands of motion clips. The model usually does not process all of them at once. Instead, it processes a small collection called a mini-batch. All pairwise relationships used by the contextual loss are computed inside that mini-batch.

This makes batch composition mathematically important. If a batch accidentally contains only one example from each class, there are no same-class pairs to pull together. If one class dominates the batch, neighborhood sizes no longer line up cleanly with class size.

\begin{conceptbox}[title=Plain-language definition of balanced sampling]
Choose the same number of examples from several classes. For example, choose four action classes and four clips from each class. The batch then contains sixteen clips.
\end{conceptbox}

\section{The two letters $P$ and $K$}

Many metric-learning implementations call this a $P\times K$ sampler:
\begin{itemize}[leftmargin=*]
  \item $P$: number of distinct classes selected for the batch;
  \item $K$: number of samples selected from each chosen class.
\end{itemize}
The total batch size is
\[
  n=PK.
\]
The contextual paper uses $k$ for the samples-per-class value, so this guide identifies $k=K$.



\section{Why batch composition is part of the method}

The paper uses a $P\times K$ sampler: choose $P$ classes and $K$ samples per class. To match the paper's notation, this guide writes
\[
  k=K,\qquad n=Pk.
\]
The contextual neighborhood size is \textbf{not a free top-$k$ tuning parameter}: the paper sets it equal to the number of samples from each class in the mini-batch.

\begin{figure}[H]
\centering
\begin{tikzpicture}[scale=0.9]
  \foreach \r/\lab/\col in {0/A/teal,1/B/gold,2/C/brick,3/D/purple}{
    \node[anchor=east] at (-0.35,-\r) {class \lab};
    \foreach \c in {0,...,3}{
      \fill[\col!70] (\c,-\r-0.32) circle (0.22);
      \node[white,font=\scriptsize\bfseries] at (\c,-\r-0.32) {\the\numexpr\c+1\relax};
    }
  }
  \draw[decorate,decoration={brace,mirror,amplitude=5pt},navy] (-0.3,-4.05)--(3.3,-4.05)
    node[midway,below=7pt] {$k=4$ samples per class};
  \draw[decorate,decoration={brace,amplitude=5pt},navy] (3.75,0.15)--(3.75,-3.6)
    node[midway,right=8pt] {$P=4$ classes};
  \node[below=18mm] at (1.5,-3.55) {$n=Pk=16$};
\end{tikzpicture}
\caption{A balanced batch is a structural assumption, not a minor dataloader choice.}
\label{fig:balanced-batch}
\end{figure}

\begin{conceptbox}[title=Why $k$ equals samples per class]
If ranking is perfect inside a balanced batch, each sample's $k$ nearest items, counting itself, are exactly the $k$ examples with the same label. This identity is what lets zero contextual loss coincide with correct within-batch ranking in the paper's proposition.
\end{conceptbox}

\section{Pairwise contrastive loss}

The contextual framework retains a standard pairwise term. In similarity form, the paper uses positive and negative margins $\delta_+$ and $\delta_-$:
\begin{equation*}
\begin{aligned}
\mathcal L_{\mathrm{contrast}}
&=
\frac{\sum_{i,j:y_{ij}=1}[\delta_+-s_{ij}]_+}
{\card{\{(i,j):y_{ij}=1,\ \delta_+-s_{ij}>0\}}}\\
&\quad+
\frac{\sum_{i,j:y_{ij}=0}[s_{ij}-\delta_-]_+}
{\card{\{(i,j):y_{ij}=0,\ s_{ij}-\delta_->0\}}}.
\end{aligned}
\paperEq{29}
\label{eq:contrastive}
\end{equation*}
Here $[a]_+=\max(a,0)$. A positive pair is penalized if its similarity is below $\delta_+$; a negative pair is penalized if its similarity is above $\delta_-$. Only margin-violating pairs contribute.

\section{Triplet and supervised contrastive baselines}

These are not part of the contextual paper's final formula, but they are critical controls for the proposed motion study.

For an anchor $a$, positive $p$, negative $n$, triplet loss is
\begin{equation}
  \mathcal L_{\mathrm{triplet}}=[D(a,p)-D(a,n)+m]_+.
  \label{eq:triplet}
\end{equation}
It asks only for a relative inequality.

For supervised contrastive learning, define $P(i)$ as the set of same-label positives for anchor $i$ and $A(i)$ as all other batch indices. With temperature $\tau$,
\begin{equation}
\mathcal L_{\mathrm{supcon}}
=\sum_i -\frac{1}{\card{P(i)}}\sum_{p\in P(i)}
\log\frac{\exp(z_i^\transpose z_p/\tau)}
{\sum_{a\in A(i)}\exp(z_i^\transpose z_a/\tau)}.
\label{eq:supcon}
\end{equation}

\begin{figure}[H]
\centering
\begin{tikzpicture}[node distance=4mm]
  \node[block,minimum width=35mm,minimum height=35mm] (pair) {};
  \node[above=1mm of pair.north] {pairwise view};
  \fill[teal] ($(pair.center)+(-0.8,0)$) circle (0.16) node[above=2pt] {$i$};
  \fill[teal] ($(pair.center)+(0.8,0.65)$) circle (0.16) node[above=2pt] {$p$};
  \fill[brick] ($(pair.center)+(0.9,-0.75)$) circle (0.16) node[below=2pt] {$n$};
  \draw[teal,very thick] ($(pair.center)+(-0.65,0.08)$)--($(pair.center)+(0.65,0.55)$);
  \draw[brick,very thick,dashed] ($(pair.center)+(-0.65,-0.08)$)--($(pair.center)+(0.72,-0.65)$);

  \node[block,right=16mm of pair,minimum width=48mm,minimum height=35mm] (ctx) {};
  \node[above=1mm of ctx.north] {contextual view};
  \coordinate (c) at (ctx.center);
  \fill[teal] ($(c)+(-1.2,0)$) circle (0.15) node[above=2pt] {$i$};
  \fill[teal] ($(c)+(1.2,0)$) circle (0.15) node[above=2pt] {$j$};
  \foreach \x/\y in {-1.6/0.7,-0.9/0.85,-0.55/-0.7,0.55/-0.7,0.9/0.85,1.6/0.7}{
    \fill[gold] ($(c)+(\x,\y)$) circle (0.11);
  }
  \draw[teal!70,thick,rounded corners] ($(c)+(-1.85,-1.05)$) rectangle ($(c)+(0.15,1.1)$);
  \draw[teal!70,thick,rounded corners] ($(c)+(-0.15,-1.05)$) rectangle ($(c)+(1.85,1.1)$);
  \node[below=2mm of ctx.south,align=center] {compare overlapping neighborhoods,\\not only one edge};
\end{tikzpicture}
\caption{Standard losses reason through pairs or triplets; contextual loss also reasons through the local graph around each sample.}
\label{fig:pair-vs-context}
\end{figure}

% ================================================================
\chapter{The Full Contextual-Similarity Pipeline}
\label{chap:pipeline}
% ================================================================

\section{One picture before the details}

\begin{figure}[H]
\centering
\resizebox{\textwidth}{!}{%
\begin{tikzpicture}[node distance=7mm]
  \node[data,minimum width=24mm] (z) {normalized\\embeddings $Z$};
  \node[block,right=of z,minimum width=23mm] (s) {similarity\\$S=ZZ^\transpose$};
  \node[block,right=of s,minimum width=23mm] (d) {distance\\$D=2-2S$};
  \node[block,right=of d,minimum width=27mm] (n) {$k+\epsilon$\\neighbor mask $N$};
  \node[block,right=of n,minimum width=28mm] (m) {shared context\\$M_+,M_-$};
  \node[block,right=of m,minimum width=25mm] (w1) {gated score\\$W_1$};
  \node[block,right=of w1,minimum width=30mm] (qe) {reciprocal query\\expansion $W_2$};
  \node[data,right=of qe,minimum width=25mm] (w) {contextual\\similarity $W$};
  \node[loss,below=11mm of w,minimum width=26mm] (lc) {$\mathcal L_{\rm context}$\\match $W$ to $Y$};
  \node[loss,below=11mm of s,minimum width=28mm] (lp) {$\mathcal L_{\rm contrast}$\\pairwise margins};
  \node[loss,below=11mm of d,minimum width=26mm] (lr) {$\mathcal L_{\rm reg}$\\space usage};
  \draw[arrow] (z)--(s); \draw[arrow] (s)--(d); \draw[arrow] (d)--(n);
  \draw[arrow] (n)--(m); \draw[arrow] (m)--(w1); \draw[arrow] (w1)--(qe); \draw[arrow] (qe)--(w);
  \draw[arrow] (w)--(lc);
  \draw[arrow] (s)--(lp);
  \draw[arrow] (s.south east) to[out=-70,in=150] (lr.north west);
\end{tikzpicture}}
\caption{The complete training computation. Every contextual quantity is ultimately a function of the cosine-similarity matrix.}
\label{fig:full-pipeline}
\end{figure}

\section{Dimensions and tensor meanings}

\begin{center}
\begin{tabularx}{\textwidth}{>{\bfseries}l l X}
\toprule
Symbol & Shape & Meaning \\
\midrule
$Z$ & $n\times d$ & one normalized embedding per mini-batch item \\
$S$ & $n\times n$ & cosine similarity, $S_{ij}=z_i^\transpose z_j$ \\
$D$ & $n\times n$ & squared Euclidean distance, $D=2\mathbf 1\mathbf 1^\transpose-2S$ entrywise \\
$Y$ & $n\times n$ & binary label agreement \\
$N$ & $n\times n$ & differentiable-in-backward binary neighborhood indicator \\
$M_+,M_-$ & $n\times n$ & shared-neighbor and shared-non-neighbor counts \\
$W_1,W_2,W$ & $n\times n$ & progressively refined contextual similarities \\
\bottomrule
\end{tabularx}
\end{center}

\begin{conceptbox}[title=The conceptual compression]
The entire paper can be remembered as
\[
\boxed{\text{embeddings}\to\text{pairwise similarities}\to\text{local graph}\to\text{graph similarity}\to\text{label graph}.}
\]
The difficult part is making the discrete local graph trainable.
\end{conceptbox}

% ================================================================
\chapter{Step 1: Differentiable Neighborhoods}
\label{chap:step1}
% ================================================================

\section{Start with a physical picture of a neighborhood}

Place every embedding as a point in space. For one anchor point $i$, sort all batch points from nearest to farthest. The first $k$ points, counting the anchor itself, form its $k$-nearest neighborhood.

\begin{figure}[H]
\centering
\begin{tikzpicture}[scale=1.05]
  \fill[teal] (0,0) circle (0.17) node[below=3pt] {anchor $i$};
  \draw[teal!70,thick,dashed] (0,0) circle (1.55);
  \foreach \x/\y/\lab in {0.6/0.4/a,-0.8/0.5/b,0.3/-1.0/c,1.2/-0.3/d,-1.4/-0.8/e,2.2/0.6/f,-2.0/1.1/g}{
    \fill[gold] (\x,\y) circle (0.12) node[above=2pt] {\scriptsize \lab};
  }
  \node[anchor=west,align=left] at (3.0,0.55) {inside boundary:\\neighbor};
  \node[anchor=west,align=left] at (3.0,-0.45) {outside boundary:\\non-neighbor};
  \draw[arrow] (2.9,0.5)--(1.1,0.2);
  \draw[arrow] (2.9,-0.4)--(2.15,0.5);
\end{tikzpicture}
\caption{A neighborhood is a binary partition induced by a distance threshold.}
\end{figure}

The paper does not merely store the list of top-$k$ indices. It expresses membership through a threshold test. That allows the forward pass to create a binary matrix and the backward pass to attach a custom gradient to the threshold.

\section{The boundary distance}

Let $p(i)$ denote the sample occupying rank $k$ in the sorted distance list for anchor $i$. Then
\[
  D(i,p(i))
\]
is the distance from anchor $i$ to the boundary item. Any sample closer than this boundary is a neighbor. A small margin $\epsilon$ slightly expands the boundary.

\begin{figure}[H]
\centering
\begin{tikzpicture}[x=1.3cm,y=1cm]
  \draw[->,slate] (0,0)--(7,0) node[right]{distance from anchor $i$};
  \foreach \x/\lab in {0/i,1/a,2/b,3/p(i),4.3/c,6/d}{
    \draw[navy,thick] (\x,-0.12)--(\x,0.12);
    \node[below=4pt] at (\x,0) {\lab};
  }
  \draw[teal,very thick] (0,0.45)--(3.5,0.45);
  \node[above,teal] at (1.75,0.45) {$N_{k+\epsilon}(i)$};
  \draw[brick,dashed,thick] (3,0.75)--(3,-0.55);
  \node[above,brick] at (3,0.75) {$D(i,p(i))$};
  \draw[gold,very thick] (3,0.7)--(3.5,0.7);
  \node[above,gold] at (3.25,0.7) {$\epsilon$};
\end{tikzpicture}
\caption{The $k$-th neighbor establishes the threshold; $\epsilon$ adds a small margin.}
\end{figure}



\section{The Heaviside operation and the heuristic gradient}

A top-$k$ neighborhood is discrete. The paper expresses the threshold with a Heaviside function:
\begin{equation*}
\text{Forward:}\quad
\theta(x)=\begin{cases}1,&x\ge 0,\\0,&x<0,\end{cases}
\qquad
\text{Backward:}\quad \frac{\partial\theta(x)}{\partial x}=\alpha,
\paperEq{1}
\label{eq:heaviside}
\end{equation*}
where $\alpha>0$ is a constant. The forward pass is exact and binary; the backward pass deliberately overrides the true derivative, which is zero almost everywhere and undefined at zero.

\begin{conceptbox}[title=Read the Heaviside rule as a switch]
The forward rule has one input number $x$ and one binary output. When $x$ is nonnegative, the switch is on and returns $1$. When $x$ is negative, the switch is off and returns $0$. The symbol $\theta$ here is a function name; it is unrelated to the neural-network parameter collection also commonly denoted by $\theta$ in other contexts.
\end{conceptbox}

The true derivative of a step function is not useful for gradient descent. Away from the jump, a small change in $x$ does not change the output, so the derivative is zero. At the jump, the derivative is not defined. The paper therefore replaces the true backward behavior with a chosen constant $\alpha>0$.

\begin{center}
\begin{tabularx}{\textwidth}{>{\bfseries}l X}
\toprule
Quantity & Interpretation \\
\midrule
$x$ & value being compared with zero \\
$\theta(x)$ & exact binary decision used in the forward computation \\
$\partial\theta/\partial x$ & sensitivity used during backpropagation \\
$\alpha$ & user-chosen positive surrogate sensitivity \\
\bottomrule
\end{tabularx}
\end{center}

\begin{warningbox}[title=What the heuristic gradient does not mean]
The paper is not claiming that the mathematical derivative of the Heaviside function is a positive constant. It deliberately defines a different backward rule so that the discrete forward computation can still provide a learning signal.
\end{warningbox}


\begin{figure}[H]
\centering
\begin{subfigure}[t]{0.46\textwidth}
\centering
\begin{tikzpicture}
\begin{axis}[
  width=\linewidth,height=5cm,axis lines=middle,xmin=-2.2,xmax=2.2,ymin=-0.15,ymax=1.25,
  xlabel={$x$},ylabel={$\theta(x)$},xtick={-2,-1,0,1,2},ytick={0,1},grid=both]
  \addplot[navy,very thick,domain=-2.2:0,samples=2] {0};
  \addplot[navy,very thick,domain=0:2.2,samples=2] {1};
  \addplot[only marks,mark=*,navy] coordinates {(0,1)};
\end{axis}
\end{tikzpicture}
\caption{Exact binary forward pass.}
\end{subfigure}\hfill
\begin{subfigure}[t]{0.46\textwidth}
\centering
\begin{tikzpicture}
\begin{axis}[
  width=\linewidth,height=5cm,axis lines=middle,xmin=-2.2,xmax=2.2,ymin=-0.15,ymax=1.25,
  xlabel={$x$},ylabel={surrogate gradient},xtick={-2,-1,0,1,2},ytick={0,1},grid=both]
  \addplot[teal,very thick,domain=-2.2:2.2,samples=2] {0.75};
  \node[teal] at (axis cs:1.45,0.92) {$\alpha$};
\end{axis}
\end{tikzpicture}
\caption{Constant backward signal.}
\end{subfigure}
\caption{The method uses an exact discrete neighborhood in the forward pass and a straight-through-style heuristic in the backward pass.}
\label{fig:heaviside}
\end{figure}

\section{The $k+\epsilon$ neighborhood}

Let $p(i)$ be the index of the $k$-th closest item to $i$, counting $i$ itself. The paper defines
\begin{equation*}
\ind_{N_{k+\epsilon}}(i,j)
=\theta\!\left(-D(i,j)+\sg(D(i,p(i)))+\epsilon\right).
\paperEq{2}
\label{eq:neighborhood}
\end{equation*}
Equivalently,
\[
  j\in N_{k+\epsilon}(i)
  \quad\Longleftrightarrow\quad
  D(i,j)\le D(i,p(i))+\epsilon.
\]

\subsection{Read Equation P2 from the inside outward}

The expression
\[
\ind_{N_{k+\epsilon}}(i,j)
=\theta\!\left(-D(i,j)+\sg(D(i,p(i)))+\epsilon\right)
\]
contains four conceptual steps.

\begin{enumerate}[leftmargin=*,itemsep=5pt]
  \item Compute the candidate distance $D(i,j)$.
  \item Compute the current boundary distance $D(i,p(i))$, where $p(i)$ is the $k$-th closest item.
  \item Compare candidate distance with boundary plus margin $\epsilon$.
  \item Convert the comparison into a binary membership value.
\end{enumerate}

\begin{center}
\begin{tabularx}{\textwidth}{>{\bfseries}l X}
\toprule
Part & Meaning \\
\midrule
$N_{k+\epsilon}(i)$ & expanded neighborhood around anchor $i$ \\
$\ind_{N_{k+\epsilon}}(i,j)$ & one if $j$ is inside that neighborhood, zero otherwise \\
$D(i,j)$ & squared Euclidean distance from anchor $i$ to candidate $j$ \\
$p(i)$ & index of the item currently at neighbor rank $k$ \\
$D(i,p(i))$ & distance to the current neighborhood boundary \\
$\epsilon$ & nonnegative margin that expands the boundary \\
$\sg(\cdot)$ & stop-gradient operation; use the value but do not differentiate through it \\
$\theta(\cdot)$ & binary threshold function \\
\bottomrule
\end{tabularx}
\end{center}

\subsection{Why the minus signs produce a distance comparison}

The Heaviside function returns one when its input is nonnegative:
\[
-D(i,j)+D(i,p(i))+\epsilon\ge 0.
\]
Move $D(i,j)$ to the other side:
\[
D(i,j)\le D(i,p(i))+\epsilon.
\]
Therefore the formula is simply a compact differentiable-in-backward way to ask whether the candidate is no farther than the current boundary plus margin.

\subsection{Numerical example}

Let $k=4$. Suppose the fourth-smallest distance from anchor $i$ is
\[
D(i,p(i))=0.70,
\]
and let $\epsilon=0.05$. The effective threshold is
\[
0.70+0.05=0.75.
\]
For candidate $j$ with $D(i,j)=0.72$,
\[
-D(i,j)+D(i,p(i))+\epsilon=-0.72+0.70+0.05=0.03\ge0,
\]
so the indicator is $1$. For candidate $r$ with distance $0.82$,
\[
-0.82+0.70+0.05=-0.07<0,
\]
so the indicator is $0$.

\begin{conceptbox}[title=The mature interpretation]
Equation P2 converts a continuous distance matrix into a directed binary adjacency matrix. The forward pass constructs an exact threshold graph. The surrogate backward rule lets gradient descent alter distances so that graph edges can appear or disappear.
\end{conceptbox}


\begin{figure}[H]
\centering
\begin{tikzpicture}[x=1.3cm,y=1cm]
  \draw[->,slate,thick] (0,0)--(7.2,0) node[right] {distance from anchor $i$};
  \foreach \x/\lab/\col in {0/i/teal,0.9/a/teal,1.6/b/teal,2.5/p/gold,3.0/c/brick,4.1/d/brick,5.5/e/brick}{
    \fill[\col] (\x,0) circle (0.12);
    \node[above=3pt] at (\x,0) {\lab};
  }
  \draw[navy,very thick] (2.5,-0.45)--(2.5,0.45);
  \node[below=7pt,navy] at (2.5,0) {$D(i,p)$};
  \draw[teal,very thick] (3.2,-0.38)--(3.2,0.38);
  \node[below=7pt,teal] at (3.2,0) {$D(i,p)+\epsilon$};
  \draw[decorate,decoration={brace,amplitude=5pt},teal] (0,0.85)--(3.2,0.85)
     node[midway,above=7pt] {$N_{k+\epsilon}(i)$};
\end{tikzpicture}
\caption{$\epsilon$ expands the boundary beyond the current $k$-th neighbor and acts as a ranking margin.}
\label{fig:epsilon-neighborhood}
\end{figure}

\section{Why stop the gradient through the boundary?}

The operator $\sg(\cdot)$ treats the $k$-th-neighbor distance as a constant during backpropagation. Otherwise, the optimization could alter the threshold itself instead of moving the candidate pair relative to the threshold. In code, this is the difference between
\[
  \code{threshold.detach()}
\]
and an ordinary differentiable threshold.

\begin{implementationbox}[title=Paper-to-code translation]
\begin{lstlisting}
D = (2.0 - 2.0 * S).clamp_min(0.0)
# topk on -D returns smallest distances
neg_dk = (-D).topk(k, dim=1).values[:, -1:]
N = GreaterThan.apply(-D + eps, neg_dk.detach())
\end{lstlisting}
The mask $N$ is binary in the forward pass. Its custom autograd function returns a constant gradient with opposite signs for the two comparison arguments.
\end{implementationbox}

\section{Why not use a sigmoid?}

A low-temperature sigmoid approximates a step, but its gradient becomes concentrated in a tiny interval around the boundary. A higher-temperature sigmoid spreads gradients but ceases to behave like a binary neighborhood. The paper reports that the exact-forward, constant-backward construction performs better in its ablations.

\begin{figure}[H]
\centering
\begin{tikzpicture}
\begin{axis}[
  width=0.78\textwidth,height=6cm,axis lines=middle,xmin=-3,xmax=3,ymin=-0.05,ymax=1.05,
  xlabel={$x$},ylabel={soft membership},grid=both,legend style={at={(0.5,-0.2)},anchor=north,legend columns=3}]
  \addplot[navy,very thick,domain=-3:0,samples=2] {0};
  \addplot[navy,very thick,domain=0:3,samples=2] {1};
  \addplot[teal,thick,domain=-3:3,samples=100] {1/(1+exp(-x/0.5))};
  \addplot[gold,thick,domain=-3:3,samples=100] {1/(1+exp(-x/0.15))};
  \legend{step,step,$\tau=0.5$,$\tau=0.15$}
\end{axis}
\end{tikzpicture}
\caption{A sharper sigmoid is a better forward approximation but has useful gradients over a narrower region.}
\label{fig:sigmoid}
\end{figure}

% ================================================================
\chapter{Step 2: Shared Neighborhood Context}
\label{chap:step2}
% ================================================================

\section{Start with a social-structure question}

Direct similarity asks whether two samples are close. Contextual similarity asks a second question: do the two samples associate with the same surrounding points?

Suppose
\[
N(i)=\{i,a,b,c\},\qquad N(j)=\{j,a,b,c\}.
\]
The two samples share three surrounding neighbors, so their local contexts are strongly aligned. By contrast, if
\[
N(r)=\{r,u,v,w\},
\]
then $i$ and $r$ have almost no local context in common.

\begin{figure}[H]
\centering
\begin{tikzpicture}[node distance=14mm]
  \node[block,minimum width=47mm,minimum height=37mm] (left) {};
  \node[above=1mm of left.north] {large neighborhood overlap};
  \fill[teal] ($(left.center)+(-1.1,0)$) circle (0.14) node[below=2pt] {$i$};
  \fill[teal] ($(left.center)+(1.1,0)$) circle (0.14) node[below=2pt] {$j$};
  \foreach \x/\y in {-0.4/0.8,0.0/-0.8,0.45/0.65}{\fill[gold] ($(left.center)+(\x,\y)$) circle (0.12);}
  \draw[teal!70,thick,rounded corners] ($(left.center)+(-1.7,-1.2)$) rectangle ($(left.center)+(0.65,1.2)$);
  \draw[teal!70,thick,rounded corners] ($(left.center)+(-0.65,-1.2)$) rectangle ($(left.center)+(1.7,1.2)$);

  \node[block,right=of left,minimum width=47mm,minimum height=37mm] (right) {};
  \node[above=1mm of right.north] {small neighborhood overlap};
  \fill[teal] ($(right.center)+(-1.2,0)$) circle (0.14) node[below=2pt] {$i$};
  \fill[brick] ($(right.center)+(1.2,0)$) circle (0.14) node[below=2pt] {$r$};
  \foreach \x/\y in {-1.55/0.75,-0.7/-0.75,0.7/0.75,1.55/-0.75}{\fill[gold] ($(right.center)+(\x,\y)$) circle (0.12);}
  \draw[teal!70,thick,rounded corners] ($(right.center)+(-1.9,-1.2)$) rectangle ($(right.center)+(-0.2,1.2)$);
  \draw[brick!70,thick,rounded corners] ($(right.center)+(0.2,-1.2)$) rectangle ($(right.center)+(1.9,1.2)$);
\end{tikzpicture}
\caption{Context compares local structures, not only the direct line between two points.}
\end{figure}

\section{Turn sets into binary rows}

For anchor $i$, create a row vector whose $p$-th entry is one when $p$ belongs to $i$'s neighborhood:
\[
  N_i=\begin{bmatrix}1&1&0&1&0\end{bmatrix}.
\]
For anchor $j$, suppose
\[
  N_j=\begin{bmatrix}1&0&0&1&1\end{bmatrix}.
\]
Multiplying matching coordinates and adding gives
\[
  N_iN_j^{\mathsf T}
  =(1)(1)+(1)(0)+(0)(0)+(1)(1)+(0)(1)=2.
\]
The result is exactly the number of positions where both rows contain one: the shared-neighbor count.



\section{Neighbor overlap as a dot product}

Let $N_{ij}=\ind_{N_{k+\epsilon}}(i,j)$. The $i$-th row $N_i$ is a binary signature of the neighborhood around $i$. The number of shared neighbors is
\begin{equation}
  M_+(i,j)=\sum_{p=1}^{n}N_{ip}N_{jp}
  =\inner{N_i}{N_j}.
  \label{eq:mplus}
\end{equation}
In matrix form,
\begin{equation*}
  M_+=NN^\transpose.
\paperEq{22}
\label{eq:mplus-matrix}
\end{equation*}
This is why the method can be implemented compactly with a matrix multiplication.

\begin{figure}[H]
\centering
\begin{tikzpicture}[scale=0.75]
  \node at (-1.1,2.6) {$N_i$};
  \foreach \c/\v in {0/1,1/1,2/0,3/1,4/0,5/0}{
    \ifnum\v=1
      \fill[teal!75] (\c,2) rectangle ++(0.8,0.8);
    \else
      \fill[mist] (\c,2) rectangle ++(0.8,0.8);
    \fi
    \draw[white] (\c,2) rectangle ++(0.8,0.8);
    \node at (\c+0.4,2.4) {\v};
  }
  \node at (-1.1,1.35) {$N_j$};
  \foreach \c/\v in {0/1,1/0,2/0,3/1,4/1,5/0}{
    \ifnum\v=1
      \fill[gold!75] (\c,0.95) rectangle ++(0.8,0.8);
    \else
      \fill[mist] (\c,0.95) rectangle ++(0.8,0.8);
    \fi
    \draw[white] (\c,0.95) rectangle ++(0.8,0.8);
    \node at (\c+0.4,1.35) {\v};
  }
  \draw[decorate,decoration={brace,mirror,amplitude=5pt},navy] (0,0.55)--(3.8,0.55)
    node[midway,below=7pt] {two shared 1s};
  \node[anchor=west] at (5.2,1.85) {$\inner{N_i}{N_j}=2$};
\end{tikzpicture}
\caption{Binary neighborhood overlap is an ordinary dot product.}
\label{fig:binary-dot}
\end{figure}

\section{Why count shared non-neighbors?}

Define $\bar N=1-N$. Then
\begin{equation}
  M_-(i,j)=\sum_{p=1}^{n}(1-N_{ip})(1-N_{jp})
  =\inner{\bar N_i}{\bar N_j}.
  \label{eq:mminus}
\end{equation}
Two neighborhood signatures can be similar because they include the same points \emph{and} exclude the same points. The $M_-$ term also prevents gradients from focusing only on entries where both neighborhood indicators are one.

\begin{figure}[H]
\centering
\begin{tikzpicture}[scale=0.92]
  \begin{scope}[xshift=-3.3cm]
    \draw[teal,very thick,fill=teal!10] (-0.7,0) circle (1.25);
    \draw[gold,very thick,fill=gold!10] (0.7,0) circle (1.25);
    \node[teal] at (-1.25,1.35) {$N(i)$};
    \node[gold] at (1.25,1.35) {$N(j)$};
    \node at (0,0) {$M_+$};
    \node[below=16mm,align=center] at (0,-1.25) {agreement on\\who is nearby};
  \end{scope}
  \begin{scope}[xshift=3.3cm]
    \draw[navy,thick,rounded corners] (-2,-1.45) rectangle (2,1.45);
    \draw[teal,very thick,fill=white] (-0.7,0) circle (0.9);
    \draw[gold,very thick,fill=white] (0.7,0) circle (0.9);
    \node at (0,-1.05) {$M_-$ region};
    \node[below=16mm,align=center] at (0,-1.25) {agreement on\\who is far away};
  \end{scope}
\end{tikzpicture}
\caption{$M_+$ measures shared inclusion; $M_-$ measures shared exclusion.}
\label{fig:mplus-mminus}
\end{figure}

\section{The intermediate contextual score $W_1$}

The paper combines the two agreements, normalizes them, and gates the result by whether $j$ is currently a neighbor of $i$:
\begin{equation*}
\boxed{
W_1(i,j)=\frac{N_{ij}}{2}
\left(
\frac{M_+(i,j)}{\sg\sum_p N_{ip}}
+
\frac{M_-(i,j)}{\sg\sum_p(1-N_{ip})}
\right).}
\paperEq{3}
\label{eq:w1}
\end{equation*}

\subsection{Read Equation P3 one layer at a time}

Rewrite the score as
\[
W_1(i,j)=
\underbrace{N_{ij}}_{\text{gate}}
\times
\underbrace{\frac{1}{2}}_{\text{average two agreements}}
\times
\left(
\underbrace{\frac{M_+(i,j)}{\sg\sum_pN_{ip}}}_{\text{shared-neighbor fraction}}
+
\underbrace{\frac{M_-(i,j)}{\sg\sum_p(1-N_{ip})}}_{\text{shared-non-neighbor fraction}}
\right).
\]

\begin{center}
\begin{tabularx}{\textwidth}{>{\bfseries}l X}
\toprule
Factor & Function \\
\midrule
$N_{ij}$ & sets the score to zero unless $j$ is currently a neighbor of $i$ \\
$M_+(i,j)$ & counts samples that are neighbors of both $i$ and $j$ \\
$\sum_pN_{ip}$ & counts the total neighbors of $i$ \\
$M_-(i,j)$ & counts samples that are non-neighbors of both $i$ and $j$ \\
$\sum_p(1-N_{ip})$ & counts the total non-neighbors of $i$ \\
$1/2$ & averages the inclusion-agreement and exclusion-agreement terms \\
$\sg$ & prevents optimization from changing denominators merely to alter normalization \\
\bottomrule
\end{tabularx}
\end{center}

\subsection{A complete numerical calculation}

Assume a batch of $n=8$ samples and $k=4$ neighbors per row. Suppose $j$ is a neighbor of $i$, so $N_{ij}=1$. Suppose the two samples share three neighbors and three non-neighbors:
\[
M_+(i,j)=3,\qquad M_-(i,j)=3.
\]
Since each row has four neighbors and four non-neighbors,
\[
\sum_pN_{ip}=4,
\qquad
\sum_p(1-N_{ip})=4.
\]
Therefore
\begin{align*}
W_1(i,j)
&=\frac{1}{2}\left(\frac{3}{4}+\frac{3}{4}\right)\\
&=\frac{1}{2}(1.5)\\
&=0.75.
\end{align*}
If $N_{ij}=0$, the outer gate makes the entire score zero even if the two binary rows agree elsewhere.

\begin{conceptbox}[title=Why this is more than pairwise similarity]
The value $W_1(i,j)$ depends on many entries of the neighborhood matrix, not only on $s_{ij}$. Changing one distance can change a neighborhood edge; that edge can change several overlaps; those overlaps can change several contextual scores.
\end{conceptbox}


The three design choices are distinct:
\begin{enumerate}[leftmargin=*,itemsep=3pt]
  \item \textbf{Gate $N_{ij}$:} contextual similarity is initially nonzero only when $j$ belongs to $i$'s neighborhood.
  \item \textbf{Average $M_+$ and $M_-$:} reward agreement about both local inclusion and exclusion.
  \item \textbf{Stop-gradient denominators:} normalize the score without training the model to manipulate neighborhood cardinality.
\end{enumerate}

\section{Algebra when $\epsilon=0$}

When $\epsilon=0$, every row has exactly $k$ ones, so $\sum_pN_{ip}=k$ and $\sum_p(1-N_{ip})=n-k$. The complement overlap simplifies:
\begin{align*}
M_-(i,j)
&=\sum_p(1-N_{ip})(1-N_{jp})\\
&=n-2k+\inner{N_i}{N_j}.
\paperEq{12}
\label{eq:mminus-simplify}
\end{align*}
Substitution gives
\begin{equation*}
W_1(i,j)=N_{ij}\left(a\inner{N_i}{N_j}+b\right),
\paperEq{14}
\label{eq:w1-simplified}
\end{equation*}
where
\begin{equation}
  a=\frac{1}{2k}+\frac{1}{2(n-k)},\qquad
  b=\frac{n-2k}{2(n-k)}.
  \label{eq:ab}
\end{equation}
Under the proposition's condition $n>2k$, both $a$ and $b$ are positive. Thus, if $N_{ij}=1$, then $W_1(i,j)>0$ even before considering how large the overlap is.

\medskip
\noindent\textbf{\textcolor{purple}{Graph interpretation.}}
Each row $N_i$ is an adjacency vector in a directed $k$-nearest-neighbor graph. The dot product $\inner{N_i}{N_j}$ counts common outgoing neighbors, so $W_1$ compares the \emph{local graph signatures} of $i$ and $j$, not merely their direct edge weight $s_{ij}$.


% ================================================================
\chapter{Step 3: Reciprocal Query Expansion and Symmetry}
\label{chap:step3}
% ================================================================

\section{Start with two refinements}

The paper adds two final refinements to the raw neighborhood-overlap score.

\begin{enumerate}[leftmargin=*,itemsep=5pt]
  \item \textbf{Reciprocity:} trust a close relation more when each point selects the other.
  \item \textbf{Query expansion:} let a query borrow evidence from its most trustworthy reciprocal neighbors.
\end{enumerate}

\section{Why nearest-neighbor relations are directional}

It is possible for $j$ to be among the nearest neighbors of $i$ while $i$ is not among the nearest neighbors of $j$. This occurs when local densities differ. A point inside a dense cluster has many very close alternatives; a point in a sparse area does not.

A reciprocal relation requires both directions:
\[
  j\in N(i)\quad\text{and}\quad i\in N(j).
\]
Binary multiplication implements logical AND:
\[
  1\cdot1=1,
\]
while every other binary product is zero.

\section{Query expansion as local evidence averaging}

Suppose $i$ has two trusted reciprocal neighbors, $a$ and $b$. Instead of using only the row $W_1(i,:)$, query expansion averages
\[
  W_1(i,:),\quad W_1(a,:),\quad W_1(b,:).
\]
A single unstable relation can therefore be moderated by consistent local evidence.

\begin{figure}[H]
\centering
\begin{tikzpicture}[node distance=8mm]
  \node[data,minimum width=25mm] (i) {query $i$};
  \node[data,above right=of i,minimum width=25mm] (a) {reciprocal\\neighbor $a$};
  \node[data,below right=of i,minimum width=25mm] (b) {reciprocal\\neighbor $b$};
  \node[block,right=22mm of i,minimum width=38mm] (avg) {average their\\contextual-score rows};
  \node[data,right=of avg,minimum width=31mm] (w2) {expanded row\\$W_2(i,:)$};
  \draw[arrow] (i)--(avg);
  \draw[arrow] (a)--(avg);
  \draw[arrow] (b)--(avg);
  \draw[arrow] (avg)--(w2);
\end{tikzpicture}
\caption{Query expansion replaces one point's estimate with a local consensus estimate.}
\end{figure}



\section{Reciprocal close neighbors}

The method recomputes a smaller $k/2+\epsilon$ neighborhood and keeps only reciprocal relationships:
\begin{equation*}
R_{k/2+\epsilon}(i,j)
=N_{k/2+\epsilon}(i,j)N_{k/2+\epsilon}(j,i).
\paperEq{4a}
\label{eq:reciprocal}
\end{equation*}
A one-sided accidental neighbor is rejected. A mutual neighbor is retained.

\begin{figure}[H]
\centering
\begin{tikzpicture}[node distance=18mm]
  \node[circle,draw=navy,thick,fill=mist,minimum size=8mm] (i) {$i$};
  \node[circle,draw=navy,thick,fill=mist,minimum size=8mm,right=of i] (j) {$j$};
  \draw[arrow,brick] (i) to[bend left=18] node[above] {$j\in N(i)$} (j);
  \node[below=8mm of $(i)!0.5!(j)$,brick] {one-sided: reject};

  \node[circle,draw=navy,thick,fill=mist,minimum size=8mm,right=55mm of j] (p) {$p$};
  \node[circle,draw=navy,thick,fill=mist,minimum size=8mm,right=of p] (q) {$q$};
  \draw[arrow,teal] (p) to[bend left=18] (q);
  \draw[arrow,teal] (q) to[bend left=18] (p);
  \node[below=8mm of $(p)!0.5!(q)$,teal] {reciprocal: retain};
\end{tikzpicture}
\caption{Reciprocity filters asymmetric nearest-neighbor accidents.}
\label{fig:reciprocity}
\end{figure}

\section{Query expansion}

The method then averages $W_1$ over the reciprocal neighbors of $i$:
\begin{equation}
  W_2(i,j)=
  \frac{\sum_p R_{k/2+\epsilon}(i,p)W_1(p,j)}
       {\sum_p R_{k/2+\epsilon}(i,p)}.
  \label{eq:w2}
\end{equation}
Finally it symmetrizes:
\begin{equation*}
  \boxed{w_{ij}=\frac12\left(W_2(i,j)+W_2(j,i)\right).}
\paperEq{4}
\label{eq:qe}
\end{equation*}

\subsection{Symbol-by-symbol interpretation of reciprocal query expansion}

The reciprocal mask is
\[
\ind_{R_{k/2+\epsilon}}(i,j)
=
\ind_{N_{k/2+\epsilon}}(i,j)
\ind_{N_{k/2+\epsilon}}(j,i).
\]
Both factors must equal one. The first says $j$ selects $i$'s local region; the second says $i$ selects $j$'s local region. Their product therefore encodes mutual selection.

The expanded score is
\[
W_2(i,j)=
\frac{\sum_p\ind_{R_{k/2+\epsilon}}(i,p)W_1(p,j)}
{\sum_p\ind_{R_{k/2+\epsilon}}(i,p)}.
\]
Read the numerator as: add the $W_1(p,j)$ scores over all reciprocal neighbors $p$ of query $i$. Read the denominator as: count how many reciprocal neighbors were included. Dividing produces an average.

\begin{center}
\begin{tabularx}{\textwidth}{>{\bfseries}l X}
\toprule
Symbol & Meaning \\
\midrule
$p$ & temporary index ranging over possible reciprocal neighbors \\
$R_{k/2+\epsilon}(i)$ & trusted mutual-neighbor set of query $i$ \\
$W_1(p,j)$ & preliminary contextual opinion of neighbor $p$ about candidate $j$ \\
$W_2(i,j)$ & local consensus opinion for query $i$ about candidate $j$ \\
$w_{ij}$ & final symmetric contextual similarity \\
\bottomrule
\end{tabularx}
\end{center}

Finally,
\[
w_{ij}=\frac{W_2(i,j)+W_2(j,i)}{2}
\]
averages both directional estimates. This guarantees $w_{ij}=w_{ji}$.


\begin{center}
\begin{tikzpicture}[node distance=8mm,scale=0.82,transform shape]
  \node[circle,draw=navy,thick,fill=sky,minimum size=8mm] (i) {$i$};
  \node[circle,draw=teal,thick,fill=lightgreen,minimum size=7mm,above right=of i] (p1) {$p_1$};
  \node[circle,draw=teal,thick,fill=lightgreen,minimum size=7mm,below right=of i] (p2) {$p_2$};
  \node[circle,draw=brick,thick,fill=rose,minimum size=7mm,right=24mm of i] (j) {$j$};
  \draw[teal,thick] (i)--(p1); \draw[teal,thick] (i)--(p2);
  \draw[dashedarrow] (p1)--(j); \draw[dashedarrow] (p2)--(j);
\end{tikzpicture}
\end{center}
\noindent\textbf{\textcolor{teal}{Statistical intuition.}}
Query expansion resembles local smoothing: a noisy pairwise estimate is replaced by an average across a small trusted neighborhood. Averaging can reduce variance, while reciprocity limits contamination from unreliable neighbors.


% ================================================================
\chapter{The Complete Objective and the Role of Every Term}
\label{chap:objective}
% ================================================================

\section{Begin with the idea of matching two matrices}

By this stage, the model has produced a contextual-similarity matrix
\[
  W=[w_{ij}],
\]
and the labels have produced a target matrix
\[
  Y=[y_{ij}].
\]
For a same-label pair, the desired target is $y_{ij}=1$. For a different-label pair, the desired target is $y_{ij}=0$. The simplest way to penalize disagreement is squared error:
\[
  (y_{ij}-w_{ij})^2.
\]

\begin{center}
\begin{tabular}{cccc}
\toprule
$y_{ij}$ & $w_{ij}$ & error & squared error \\
\midrule
1 & 0.9 & $0.1$ & $0.01$ \\
1 & 0.2 & $0.8$ & $0.64$ \\
0 & 0.1 & $-0.1$ & $0.01$ \\
0 & 0.8 & $-0.8$ & $0.64$ \\
\bottomrule
\end{tabular}
\end{center}

Averaging this quantity across all distinct pairs produces the contextual loss.

\section{Why the final objective contains three terms}

The final loss combines three pressures:
\begin{itemize}[leftmargin=*]
  \item contextual structure should match the label structure;
  \item direct cosine similarities should remain locally meaningful across batches;
  \item the whole embedding space should not collapse into a narrow angular region.
\end{itemize}
No single term performs all three roles equally well.



\section{Contextual loss}

The final contextual matrix $W=[w_{ij}]$ is trained to match the label-similarity matrix $Y$:
\begin{equation*}
  \boxed{\mathcal L_{\mathrm{context}}
  =\frac{1}{n^2}\sum_{i\ne j}(y_{ij}-w_{ij})^2.}
\paperEq{5}
\label{eq:lcontext}
\end{equation*}

\subsection{Read the contextual loss as four nested operations}

The formula
\[
\mathcal L_{\mathrm{context}}
=\frac{1}{n^2}\sum_{i\ne j}(y_{ij}-w_{ij})^2
\]
performs four operations:
\begin{enumerate}[leftmargin=*,itemsep=4pt]
  \item compare target $y_{ij}$ with prediction $w_{ij}$;
  \item square the error;
  \item add errors over all distinct pairs;
  \item divide by $n^2$ to keep the scale approximately stable as batch size changes.
\end{enumerate}

\begin{center}
\begin{tabularx}{\textwidth}{>{\bfseries}l X}
\toprule
Part & Meaning \\
\midrule
$\mathcal L_{\mathrm{context}}$ & one nonnegative scalar loss for the entire mini-batch \\
$n$ & number of batch samples \\
$i\ne j$ & exclude trivial self-pairs \\
$y_{ij}$ & binary target relation \\
$w_{ij}$ & learned contextual-similarity prediction in $[0,1]$ \\
$(y_{ij}-w_{ij})^2$ & squared disagreement for one pair \\
\bottomrule
\end{tabularx}
\end{center}

The loss equals zero only when every included squared term equals zero. Since a square is zero exactly when its argument is zero, this requires
\[
  w_{ij}=y_{ij}\qquad\text{for every }i\ne j.
\]
This observation is central to the ranking proposition.

The diagonal is excluded because every sample is trivially identical to itself.

The derivative with respect to a contextual score is ordinary MSE:
\begin{equation}
  \frac{\partial\mathcal L_{\mathrm{context}}}{\partial w_{ij}}
  =\frac{2}{n^2}(w_{ij}-y_{ij}).
  \label{eq:msegrad}
\end{equation}
The unusual part is the chain from $w_{ij}$ back through query expansion, neighborhood intersections, and discrete threshold comparisons to the embedding similarities.

\section{Similarity regularizer}

The paper also regulates the mean pairwise cosine similarity:
\begin{equation*}
  \mathcal L_{\mathrm{reg}}
  =\left(\tilde s-\frac{1}{n^2}\sum_{i,j}s_{ij}\right)^2.
\paperEq{7}
\label{eq:lreg}
\end{equation*}
If all embeddings crowd into a small angular region, their mean similarity becomes too high. This term encourages broader use of the unit sphere.

\section{The full objective}

\begin{equation*}
  \boxed{
  \mathcal L_{\mathrm{ours}}
  =\lambda\mathcal L_{\mathrm{context}}
  +(1-\lambda)\mathcal L_{\mathrm{contrast}}
  +\gamma\mathcal L_{\mathrm{reg}}.}
\paperEq{6}
\label{eq:full-objective}
\end{equation*}

\subsection{How the weighting coefficients work}

The coefficients $\lambda$ and $\gamma$ do not represent probabilities. They determine relative gradient scale.

\begin{itemize}[leftmargin=*]
  \item $\lambda=1$ removes the contrastive component and uses contextual loss plus regularization.
  \item $\lambda=0$ removes contextual loss and uses contrastive loss plus regularization.
  \item an intermediate $\lambda$ combines both signals.
  \item $\gamma=0$ disables the mean-similarity regularizer.
\end{itemize}

The coefficient $(1-\lambda)$ is a design choice that makes the two principal weights sum to one. It does not imply that the two unweighted losses naturally have equal numerical scale.


\begin{figure}[H]
\centering
\begin{tikzpicture}[node distance=10mm]
  \node[loss,minimum width=42mm] (ctx) {$\mathcal L_{\rm context}$\\correct local ranking graph};
  \node[loss,right=of ctx,minimum width=42mm] (con) {$\mathcal L_{\rm contrast}$\\maintain pairwise margins};
  \node[loss,right=of con,minimum width=42mm] (reg) {$\mathcal L_{\rm reg}$\\avoid angular crowding};
  \node[block,below=15mm of con,minimum width=42mm] (sum) {weighted sum\\$\mathcal L_{\rm ours}$};
  \draw[arrow] (ctx)--(sum); \draw[arrow] (con)--(sum); \draw[arrow] (reg)--(sum);
\end{tikzpicture}
\caption{The three terms solve different failure modes.}
\label{fig:loss-roles}
\end{figure}

\section{The decomposability gap}

A batch can be perfectly ranked even while the entire dataset is not. Once $\mathcal L_{\mathrm{context}}=0$ for an easy batch, that batch provides no further signal. The pairwise contrastive term continues enforcing absolute margins and helps bridge this difference between batch-wise and dataset-wise ranking.

\begin{figure}[H]
\centering
\begin{tikzpicture}[node distance=8mm]
  \node[block,minimum width=42mm,minimum height=24mm] (batch) {sampled batch\\all positives ahead of negatives\\$\mathcal L_{\rm context}=0$};
  \node[block,right=25mm of batch,minimum width=42mm,minimum height=24mm] (data) {full dataset\\some unseen hard negatives\\ranking still imperfect};
  \draw[dashedarrow,brick] (batch)-- node[above,align=center] {decomposability\\gap} (data);
  \node[loss,below=12mm of $(batch)!0.5!(data)$] (c) {contrastive margins keep acting};
  \draw[arrow] (c)--(batch); \draw[arrow] (c)--(data);
\end{tikzpicture}
\caption{Why the hybrid objective is stronger than contextual loss alone.}
\label{fig:decomp-gap}
\end{figure}

\section{Hyperparameter map}

\begin{longtable}{>{\bfseries}p{0.12\textwidth}p{0.22\textwidth}p{0.49\textwidth}}
\toprule
Parameter & Mathematical role & Practical interpretation \\
\midrule
$k$ & neighbor count & must equal samples per class in the balanced batch under the paper's theory \\
$\epsilon$ & expands the $k$-th-neighbor threshold & a flexible ranking margin in squared-distance units \\
$\alpha$ & surrogate derivative of $\theta$ & scales gradients through comparisons; largely redundant with learning rate \\
$\delta_+$ & positive similarity margin & positives below it are pulled closer \\
$\delta_-$ & negative similarity margin & negatives above it are pushed apart \\
$\lambda$ & contextual--contrastive mixture & $1$ means contextual only; $0$ means contrastive only \\
$\gamma$ & similarity-regularizer weight & controls pressure to use more of the sphere \\
$\tilde s$ & desired mean cosine similarity & target average angular spread \\
\bottomrule
\end{longtable}

% ================================================================
\chapter{Why Zero Contextual Loss Means Correct Ranking}
\label{chap:proof}
% ================================================================

\section{How to read a mathematical proposition}

A proposition has assumptions and a conclusion. The assumptions describe the setting in which the statement is guaranteed. Removing an assumption can make the conclusion false.

The phrase ``if and only if,'' written $\Longleftrightarrow$, contains two directions:
\[
  A\Longrightarrow B
  \qquad\text{and}\qquad
  B\Longrightarrow A.
\]
For this paper:
\begin{itemize}[leftmargin=*]
  \item direction 1 asks whether a correct ranking forces contextual loss to zero;
  \item direction 2 asks whether zero contextual loss forces the ranking to be correct.
\end{itemize}
The second direction is usually more informative because it rules out an incorrect ranking that somehow achieves a perfect loss value.

\section{What ``correct ranking'' means formally}

For anchor $i$, let $p$ be any positive sample with $y_{ip}=1$, and let $j$ be any negative sample with $y_{ij}=0$. Correct ranking means
\[
  s_{ip}>s_{ij}.
\]
Read this as: every same-class sample must have larger cosine similarity to the anchor than every different-class sample.

Because distance decreases when similarity increases, the equivalent distance statement is
\[
  D(i,p)<D(i,j).
\]



\section{The paper's proposition}
\label{sec:proposition}

\begin{proposition}[Liao--Tsiligkaridis--Kulis, Proposition 4.1]
Consider a balanced batch of size $n$ with exactly $k$ samples from each class, $n$ divisible by $k$, $n>2k$, $k\ge2$, and $\epsilon=0$. Then
\[
\mathcal L_{\mathrm{context}}=0
\]
if and only if every positive sample is ranked ahead of every negative sample for every anchor:
\[
  s_{ip}>s_{ij}
  \quad\text{whenever}\quad y_{ip}=1,\ y_{ij}=0.
\]
\end{proposition}

\subsection{Translate every quantifier in the ranking condition}

The paper writes a condition of the form
\[
  s_{ip}>s_{ij},
  \qquad
  \forall p,j\text{ such that }y_{ip}=1\text{ and }y_{ij}=0,
  \qquad
  \forall i.
\]
Read the quantifier $\forall$ as ``for every.'' The complete sentence is:

\begin{quote}
For every anchor sample $i$, for every same-label sample $p$, and for every different-label sample $j$, the similarity between $i$ and $p$ must be greater than the similarity between $i$ and $j$.
\end{quote}

This is stronger than requiring the nearest neighbor to be correct. It requires all positives to appear ahead of all negatives inside the batch.

\begin{conceptbox}[title=Why counting the anchor itself matters]
Each class contributes exactly $k$ samples to the batch, including the anchor. Under perfect ranking, the first $k$ positions are precisely the complete same-class group. This is why the neighborhood size and samples-per-class count must match in the proposition.
\end{conceptbox}


\begin{conceptbox}[title=Plain-language translation]
Under the balanced-batch assumptions, the contextual loss has the right global minimum: it is zero exactly when every anchor's same-class examples occupy the full top-$k$ neighborhood and every different-class example lies below them.
\end{conceptbox}

\section{Forward direction: correct ranking implies zero loss}

Assume every sample is correctly ranked.

\begin{enumerate}[leftmargin=*,label=\textbf{Step \arabic*.},itemsep=5pt]
  \item Each class contributes exactly $k$ examples. Therefore the top-$k$ neighborhood of anchor $i$ is exactly its class:
  \[
    N_{ij}=y_{ij}.
  \]
  \item If $N_{ij}=1$, then $i$ and $j$ have identical neighborhoods of size $k$, so
  \[
    \inner{N_i}{N_j}=k.
  \]
  \item Substituting into the simplified $W_1$ formula yields $W_1(i,j)=N_{ij}=y_{ij}$.
  \item Reciprocal query expansion averages identical rows inside the same class, so it leaves these values unchanged.
  \item Hence $w_{ij}=y_{ij}$ for all off-diagonal pairs and the MSE is zero.
\end{enumerate}

\section{Backward direction: zero loss implies correct ranking}

Assume $\mathcal L_{\mathrm{context}}=0$ but ranking is not correct. Because the batch contains exactly $k$ positives for each anchor, including itself, a missing positive from the top-$k$ set forces at least one negative intruder into the top-$k$ set. So there exists a negative pair $(i,j)$ with
\[
  y_{ij}=0,\qquad N_{ij}=1.
\]
From \cref{eq:w1-simplified}, $a>0$ and $b>0$, hence
\[
  W_1(i,j)=a\inner{N_i}{N_j}+b>0.
\]
Query expansion is an average of nonnegative quantities with positive weights, so the corresponding contextual value remains positive:
\[
  w_{ij}>0.
\]
But $y_{ij}=0$, so $(y_{ij}-w_{ij})^2>0$, contradicting zero loss. Therefore ranking must be correct.

\begin{figure}[H]
\centering
\begin{tikzpicture}[scale=0.58]
  \foreach \r in {0,...,11}{
    \foreach \c in {0,...,11}{
      \pgfmathtruncatemacro{\same}{(int(\r/4)==int(\c/4))?1:0}
      \ifnum\same=1\fill[teal!75] (\c,-\r) rectangle ++(1,-1);\else\fill[mist] (\c,-\r) rectangle ++(1,-1);\fi
      \draw[white,line width=0.25pt] (\c,-\r) rectangle ++(1,-1);
    }
  }
  \draw[navy,thick] (0,0) rectangle (12,-12);
  \node[above=5mm] at (6,0) {perfect-neighborhood matrix $N=Y$};
  \draw[brick,very thick] (3,-0.05) rectangle (5,-1.95);
  \node[right=5mm,brick,align=left] at (12,-1) {a negative intruder creates\\a nonzero off-block entry};
\end{tikzpicture}
\caption{At the optimum, the neighbor graph has exactly the same block structure as label agreement.}
\label{fig:proof-matrix}
\end{figure}

\section{What the assumptions are doing}

\begin{center}
\small
\begin{tabularx}{\textwidth}{>{\bfseries}p{0.31\textwidth} X}
\toprule
Assumption & Why it matters \\
\midrule
balanced $k$ samples per class & makes the desired top-$k$ set exactly equal to the same-label set \\
$\epsilon=0$ & fixes each neighborhood cardinality at exactly $k$ \\
$n>2k$ & ensures the constant $b=(n-2k)/(2(n-k))$ is positive in the proof \\
self included & makes a class of $k$ examples correspond to a neighborhood of size $k$ \\
outside these assumptions & the loss may remain useful empirically, but the proposition's zero-loss/correct-ranking equivalence is not guaranteed \\
\bottomrule
\end{tabularx}
\end{center}

% ================================================================
\chapter{Gradient Intuition and Semantic Consistency}
\label{chap:gradients}
% ================================================================

\section{Begin with the difference between a forward value and a backward signal}

The forward pass answers: what score and loss do the current embeddings produce? The backward pass answers: in which direction should the embeddings and parameters change?

These are separate questions. A function may produce a sensible forward value but have a derivative that is unusable for learning. The hard indicator is the main example: its forward output correctly represents membership, but its true gradient is zero almost everywhere.

\section{What a pairwise gradient means geometrically}

Consider a similarity $s_{ij}$. If
\[
  \frac{\partial\mathcal L}{\partial s_{ij}}<0,
\]
then increasing $s_{ij}$ would decrease the loss, so gradient descent tends to make the pair more similar. If
\[
  \frac{\partial\mathcal L}{\partial s_{ij}}>0,
\]
then decreasing $s_{ij}$ would decrease the loss, so gradient descent tends to separate the pair.

For contextual loss, this derivative is influenced by many neighborhood relations. A pair can receive a gradient not only because of its own label but because changing that pair changes neighborhood membership and therefore changes contextual scores for several other pairs.



\section{Simplified gradient view}

For $\epsilon=0$, the paper writes
\begin{equation*}
  W_1(i,j)=N_{ij}\left(a\inner{N_i}{N_j}+b\right).
\paperEq{8}
\label{eq:paper8}
\end{equation*}
Consider a negative intruder $j\in N(i)$, so $y_{ij}=0$ but $N_{ij}=1$. The MSE wants to reduce $W_1(i,j)$. It can do so in two ways:

\begin{enumerate}[leftmargin=*,itemsep=4pt]
  \item remove the direct neighborhood edge $i\to j$;
  \item reduce overlap between the neighborhood signatures $N_i$ and $N_j$.
\end{enumerate}
The second route changes distances between $i$ and members of $j$'s neighborhood. That is the key distinction from a purely pairwise loss.

\begin{figure}[H]
\centering
\begin{tikzpicture}[scale=1.0]
  \node[circle,draw=navy,thick,fill=sky,minimum size=9mm] (i) at (-2.2,0) {$i$};
  \node[circle,draw=brick,thick,fill=rose,minimum size=9mm] (j) at (2.2,0) {$j$};
  \foreach \x/\y/\name in {1.35/1.1/p1,2.7/1.0/p2,1.25/-1.05/p3,2.8/-1.0/p4}{
    \node[circle,draw=gold,fill=amber,minimum size=6mm] (\name) at (\x,\y) {};
    \draw[gold,thick] (j)--(\name);
  }
  \draw[brick,very thick] (i)-- node[above] {wrong negative edge} (j);
  \draw[brick,dashed,thick,-{Latex}] (i) to[bend left=18] (p1);
  \draw[brick,dashed,thick,-{Latex}] (i) to[bend right=18] (p3);
  \node[below=16mm,align=center] at (0,0) {contextual gradient separates $i$ from\\the \emph{context} around $j$, not only from $j$};
\end{tikzpicture}
\caption{A contextual error propagates through a neighborhood.}
\label{fig:context-gradient}
\end{figure}

\section{Why apparently ``wrong'' pairwise gradients can generalize}

The paper observes that a contextual gradient may sometimes pull apart same-label examples or push together different-label examples. Pairwise supervision would call that wrong. Contextual supervision asks a different question: is the binary label relation consistent with the local semantic structure?

The paper's gradient-clamping experiment enforces the pairwise label signs
\begin{equation*}
\partial D(i,j)\leftarrow
\begin{cases}
\max(\partial D(i,j),0),&y_{ij}=1,\\
\min(\partial D(i,j),0),&y_{ij}=0,
\end{cases}
\paperEq{10}
\label{eq:clamp}
\end{equation*}
and reports higher train retrieval but lower test retrieval. Its interpretation is that some label-inconsistent gradients act as regularization against semantically inconsistent labels.

\begin{warningbox}[title=Critical distinction for the pose proposal]
This argument concerns noisy or semantically incomplete \emph{labels}. Pose-estimation corruption is noisy \emph{input}. A corrupted skeleton may move an entire local neighborhood coherently or destroy it. Therefore the claim ``context helps corrupted poses'' remains an empirical hypothesis.
\end{warningbox}

\section{Stepwise gradient interpretation}

The supplementary analysis can be summarized as follows:

\begin{center}
\begin{tabularx}{\textwidth}{>{\bfseries}p{0.17\textwidth}X X}
\toprule
Component & Positive pair error & Negative pair error \\
\midrule
Step 1: $N_{ij}$ & positive outside top-$k$ is pulled inward & negative inside top-$k$ is pushed outward \\
Step 2: overlap & align the two local contexts & separate each sample from the other's local context \\
Step 3: query expansion & smooth toward trusted-neighbor consensus & suppress one-sided accidental relations \\
\bottomrule
\end{tabularx}
\end{center}

% ================================================================
\chapter{A Worked Numerical Example}
\label{chap:worked}
% ================================================================

\section{Setup}

Use a balanced batch with $n=12$, three classes, and $k=4$ samples per class:
\[
A=\{1,2,3,4\},\quad B=\{5,6,7,8\},\quad C=\{9,10,11,12\}.
\]
Suppose anchor $1$ has a damaged neighborhood
\[
N(1)=\{1,2,3,5\},
\]
so positive $4$ is missing and negative $5$ is an intruder. Let
\[
N(2)=\{2,1,3,4\},\qquad N(5)=\{5,6,7,1\}.
\]

\section{A true positive pair: $(1,2)$}

The shared neighbors are
\[
N(1)\cap N(2)=\{1,2,3\},\qquad M_+(1,2)=3.
\]
The union contains five items, so seven of the twelve items are excluded by both neighborhoods:
\[
M_-(1,2)=12-\card{N(1)\cup N(2)}=12-5=7.
\]
Because $2\in N(1)$, the gate is one. Therefore
\begin{align}
W_1(1,2)
&=\frac12\left(\frac34+\frac78\right)\\
&=\frac{13}{16}=0.8125.
\end{align}
The target is $y_{12}=1$, so this pair contributes
\[
(1-0.8125)^2\approx0.0352
\]
before the global $1/n^2$ factor.

\section{A wrong negative pair: $(1,5)$}

Here
\[
N(1)\cap N(5)=\{1,5\},\qquad M_+(1,5)=2.
\]
The union has six items, so $M_-(1,5)=6$. Since $5\in N(1)$,
\begin{align}
W_1(1,5)
&=\frac12\left(\frac24+\frac68\right)\\
&=0.625.
\end{align}
But $y_{15}=0$, so the contribution is
\[
(0-0.625)^2=0.390625.
\]
This is a strong signal to remove the intruder and separate the surrounding contexts.

\section{A missing positive pair: $(1,4)$}

Because $4\notin N(1)$, the gate makes
\[
W_1(1,4)=0
\]
even though $y_{14}=1$. Its squared error is $1$, the largest possible. Step 1 therefore pulls sample $4$ toward the top-$k$ boundary of anchor $1$.

\begin{figure}[H]
\centering
\begin{tikzpicture}[node distance=7mm]
  \node[circle,draw=navy,thick,fill=sky,minimum size=9mm] (a1) {$1_A$};
  \node[circle,draw=teal,fill=lightgreen,minimum size=7mm,above right=of a1] (a2) {$2_A$};
  \node[circle,draw=teal,fill=lightgreen,minimum size=7mm,right=of a1] (a3) {$3_A$};
  \node[circle,draw=teal,fill=white,minimum size=7mm,below right=of a1] (a4) {$4_A$};
  \node[circle,draw=brick,fill=rose,minimum size=7mm,above=17mm of a3] (b5) {$5_B$};
  \draw[teal,thick] (a1)--(a2); \draw[teal,thick] (a1)--(a3);
  \draw[brick,very thick] (a1)--(b5);
  \draw[teal,dashed,very thick,-{Latex}] (a1)--(a4);
  \node[right=10mm of b5,brick] {intruder: push out};
  \node[right=10mm of a4,teal] {missing positive: pull in};
\end{tikzpicture}
\caption{The numerical example as a neighborhood-repair problem.}
\label{fig:worked-example}
\end{figure}

\section{What query expansion would do}

Suppose $1$ and $2$ are reciprocal close neighbors. Then $W_2(1,j)$ averages the row of $W_1$ for anchors $1$ and $2$. Since anchor $2$ has the correct neighbor $4$ and excludes intruder $5$, its row supplies a corrective vote: $W_2(1,4)$ rises and $W_2(1,5)$ falls. This is the local-consensus interpretation of query expansion.

% ================================================================
\chapter{Adapting the Mathematics to Pose Sequences}
\label{chap:pose}
% ================================================================

\section{The input tensor}

A pose sequence can be written
\begin{equation}
  x\in\R^{T\times J\times C},
  \label{eq:pose-tensor}
\end{equation}
where $T$ is frames, $J$ is joints, and $C\in\{2,3\}$ is coordinate dimension. A sequence encoder produces one vector:
\begin{equation}
  h=f_\theta(x,m),\qquad
  z=\frac{g_\phi(h)}{\norm{g_\phi(h)}_2}.
  \label{eq:pose-encoder}
\end{equation}
Here $m$ can be a validity mask for missing joints or frames, $f_\theta$ is MotionBERT or another temporal encoder, and $g_\phi$ is a small projection head.

\begin{figure}[H]
\centering
\resizebox{0.94\textwidth}{!}{%
\begin{tikzpicture}[node distance=9mm]
  \node[data,minimum width=30mm] (pose) {pose tensor\\$T\times J\times C$};
  \node[block,right=of pose,minimum width=31mm] (norm) {root/scale\\normalization};
  \node[block,right=of norm,minimum width=35mm] (enc) {temporal encoder\\MotionBERT};
  \node[block,right=of enc,minimum width=30mm] (pool) {masked pooling\\or CLS token};
  \node[block,right=of pool,minimum width=28mm] (head) {projection head\\$g_\phi$};
  \node[data,right=of head,minimum width=26mm] (z) {unit motion\\embedding $z$};
  \draw[arrow] (pose)--(norm); \draw[arrow] (norm)--(enc); \draw[arrow] (enc)--(pool); \draw[arrow] (pool)--(head); \draw[arrow] (head)--(z);
\end{tikzpicture}}
\caption{The loss does not care whether $z$ came from an image or a motion encoder. Only the encoder-side computation changes.}
\label{fig:pose-encoder}
\end{figure}

\section{Pose normalization}

For joint $j$ at frame $t$, let $r$ be a root joint and $b$ a body-scale statistic such as mean bone length. A simple normalization is
\begin{equation}
  \tilde x_{t,j}=\frac{x_{t,j}-x_{t,r}}{b+\eta},
  \label{eq:pose-normalize}
\end{equation}
with small $\eta>0$ for stability. This removes global translation and much of body-size variation. View normalization, if used, should be treated as a separate design choice because it can remove information or create artifacts.

\section{Controlled corruption operators}

The main experiment corrupts the query and leaves the gallery clean. Let corruption severity be $c$.

\subsection{Gaussian joint jitter}
\begin{equation}
  \tilde x_{t,j}=x_{t,j}+\sigma(c)\xi_{t,j},
  \qquad \xi_{t,j}\sim\mathcal N(0,I_C).
  \label{eq:jitter}
\end{equation}
Use scale-normalized coordinates so $\sigma$ has comparable meaning across subjects.

\subsection{Joint or limb dropout}
Let $m_{t,j}\sim\operatorname{Bernoulli}(1-p(c))$. Then
\begin{equation}
  \tilde x_{t,j}=m_{t,j}x_{t,j},
  \label{eq:jointdrop}
\end{equation}
while a companion mask tells the encoder which zeros are missing observations. Limb dropout uses a correlated mask over an anatomically connected subset.

\subsection{Frame dropout}
With frame mask $u_t\sim\operatorname{Bernoulli}(1-q(c))$,
\begin{equation}
  \tilde x_t=u_tx_t,
  \label{eq:framedrop}
\end{equation}
or remove the frame and resample the remaining sequence to length $T$.

\subsection{View rotation in 3D}
For a yaw angle $\psi(c)$,
\begin{equation}
\tilde x_{t,j}=R_y(\psi)x_{t,j},\qquad
R_y(\psi)=
\begin{bmatrix}
\cos\psi&0&\sin\psi\\
0&1&0\\
-\sin\psi&0&\cos\psi
\end{bmatrix}.
\label{eq:rotation}
\end{equation}

\begin{figure}[H]
\centering
\begin{tikzpicture}[scale=0.88]
% clean skeleton
\begin{scope}[xshift=-5.2cm]
  \node[above] at (0,2.5) {clean};
  \draw[navy,thick] (0,2)--(0,0.5)--(-0.7,-0.6); \draw[navy,thick] (0,0.5)--(0.7,-0.6);
  \draw[navy,thick] (0,1.45)--(-0.8,0.9); \draw[navy,thick] (0,1.45)--(0.8,0.9);
  \fill[teal] (0,2.2) circle (0.2);
  \foreach \x/\y in {0/2,0/1.45,0/0.5,-0.8/0.9,0.8/0.9,-0.7/-0.6,0.7/-0.6}{\fill[navy] (\x,\y) circle (0.08);}
\end{scope}
% jitter
\begin{scope}[xshift=-1.75cm]
  \node[above] at (0,2.5) {jitter};
  \draw[navy,thick] (0.08,1.95)--(-0.06,0.55)--(-0.55,-0.75); \draw[navy,thick] (-0.06,0.55)--(0.84,-0.45);
  \draw[navy,thick] (0.08,1.4)--(-0.95,0.78); \draw[navy,thick] (0.08,1.4)--(0.68,1.03);
  \fill[teal] (-0.05,2.2) circle (0.2);
  \foreach \x/\y in {0.08/1.95,-0.06/1.4,-0.06/0.55,-0.95/0.78,0.68/1.03,-0.55/-0.75,0.84/-0.45}{\fill[navy] (\x,\y) circle (0.08);}
\end{scope}
% joint missing
\begin{scope}[xshift=1.75cm]
  \node[above] at (0,2.5) {missing limb};
  \draw[navy,thick] (0,2)--(0,0.5)--(-0.7,-0.6); \draw[navy,thick] (0,0.5)--(0.7,-0.6);
  \draw[navy,thick] (0,1.45)--(-0.8,0.9);
  \draw[brick,dashed,thick] (0,1.45)--(0.8,0.9);
  \fill[teal] (0,2.2) circle (0.2);
  \draw[brick,thick] (0.68,0.78)--(0.92,1.02); \draw[brick,thick] (0.92,0.78)--(0.68,1.02);
\end{scope}
% dropped frames
\begin{scope}[xshift=5.2cm]
  \node[above] at (0,2.5) {dropped frames};
  \foreach \x in {-0.9,0,0.9}{\draw[navy,thick,opacity=0.8] (\x,2)--(\x,1)--(\x-0.35,0.35);}
  \draw[brick,very thick] (-0.25,0.25)--(0.25,2.2); \draw[brick,very thick] (0.25,0.25)--(-0.25,2.2);
\end{scope}
\end{tikzpicture}
\caption{Schematic pose corruptions. These are controlled test interventions, not claims that one synthetic operator fully represents a real pose estimator.}
\label{fig:pose-corruptions}
\end{figure}

\section{The central hypothesis in mathematical form}

Let $M_\ell(c)$ be a retrieval metric for loss $\ell$ at corruption severity $c$. Define degradation
\begin{equation}
  \Delta_\ell(c)=M_\ell(0)-M_\ell(c).
  \label{eq:degradation}
\end{equation}
The proposed hypothesis is not merely that contextual loss has larger clean $M(0)$, but that
\begin{equation}
  \Delta_{\mathrm{context}}(c)
  <\Delta_{\mathrm{baseline}}(c)
  \quad\text{over a meaningful severity range.}
  \label{eq:hypothesis}
\end{equation}

% ================================================================
\chapter{Implementation Blueprint}
\label{chap:implementation}
% ================================================================

\section{Recommended staged implementation}

\begin{figure}[H]
\centering
\begin{tikzpicture}[node distance=5mm]
  \node[impl,minimum width=42mm] (s0) {Stage 0\\split manifest, evaluator, corruptions};
  \node[impl,right=of s0,minimum width=42mm] (s1) {Stage 1\\frozen MotionBERT features};
  \node[impl,right=of s1,minimum width=42mm] (s2) {Stage 2\\train projection heads};
  \node[impl,below=10mm of s1,minimum width=42mm] (s3) {Stage 3\\fine-tune full encoder};
  \node[impl,right=of s3,minimum width=42mm] (s4) {Stage 4\\real estimator shift / sports};
  \draw[arrow] (s0)--(s1); \draw[arrow] (s1)--(s2); \draw[arrow] (s2) to[bend left=20] (s3); \draw[arrow] (s3)--(s4);
\end{tikzpicture}
\caption{A low-risk sequence that isolates bugs and scientific effects.}
\label{fig:stages}
\end{figure}

\begin{enumerate}[leftmargin=*,itemsep=5pt]
  \item \textbf{Freeze the encoder first.} Validate query/gallery construction, cosine ranking, metrics, and deterministic corruption curves.
  \item \textbf{Train only a projection head.} Compare losses while holding the representation backbone fixed.
  \item \textbf{Fine-tune the encoder only after stability.} This prevents an encoder bug or optimizer difference from masquerading as a loss effect.
  \item \textbf{Add a corruption-aware baseline.} Compare ordinary SupCon with and without the same corruption augmentation, then add a MaskCLR-style control if feasible.
\end{enumerate}

\section{Reference PyTorch implementation}

The following implementation follows the paper's equations, parameterizes $k$, uses exact forward masks with custom backward gradients, and handles empty contrastive sets safely.

\begin{lstlisting}
from __future__ import annotations

from dataclasses import dataclass
from typing import Dict, Tuple

import torch
import torch.nn as nn
import torch.nn.functional as F


class GreaterThanSTE(torch.autograd.Function):
    """Exact >= in forward; constant surrogate gradient in backward."""

    @staticmethod
    def forward(ctx, x: torch.Tensor, y: torch.Tensor, alpha: float):
        ctx.alpha = float(alpha)
        return (x >= y).to(dtype=x.dtype)

    @staticmethod
    def backward(ctx, grad_output: torch.Tensor):
        a = ctx.alpha
        return grad_output * a, -grad_output * a, None


def contextual_similarity(
    embeddings: torch.Tensor,
    *,
    k: int,
    eps: float,
    alpha: float = 10.0,
) -> Tuple[torch.Tensor, Dict[str, torch.Tensor]]:
    """Compute W from normalized or unnormalized [n, d] embeddings."""
    if embeddings.ndim != 2:
        raise ValueError("embeddings must have shape [batch, dim]")
    n = embeddings.shape[0]
    if k < 2 or k > n:
        raise ValueError(f"k must be in [2, {n}], got {k}")
    if k % 2 != 0:
        raise ValueError("k must be even for the k/2 reciprocal step")

    z = F.normalize(embeddings, p=2, dim=1)
    S = z @ z.T
    D = (2.0 - 2.0 * S).clamp_min(0.0)

    # Step 1: k+eps neighborhood.
    neg_topk = (-D).topk(k, dim=1).values
    neg_dk = neg_topk[:, -1:]
    N = GreaterThanSTE.apply(-D + eps, neg_dk.detach(), alpha)

    # Step 2: shared neighbors and shared non-neighbors.
    N_count = N.sum(dim=1, keepdim=True).detach().clamp_min(1.0)
    M_plus = (N @ N.T) / N_count

    N_bar = 1.0 - N
    N_bar_count = N_bar.sum(dim=1, keepdim=True).detach().clamp_min(1.0)
    M_minus = (N_bar @ N_bar.T) / N_bar_count

    W1 = 0.5 * (M_plus + M_minus) * N

    # Step 3: reciprocal k/2 neighborhood and query expansion.
    k2 = k // 2
    neg_dk2 = neg_topk[:, k2 - 1 : k2]
    N2 = GreaterThanSTE.apply(-D + eps, neg_dk2.detach(), alpha)
    R = N2 * N2.T
    R_count = R.sum(dim=1, keepdim=True).clamp_min(1.0)
    W2 = (R @ W1) / R_count
    W = 0.5 * (W2 + W2.T)

    aux = {"z": z, "S": S, "D": D, "N": N, "W1": W1, "R": R}
    return W, aux


@dataclass(frozen=True)
class ContextualLossConfig:
    k: int
    eps: float = 0.05
    alpha: float = 10.0
    lam: float = 0.2
    gamma: float = 0.1
    target_mean_similarity: float = 0.0
    positive_margin: float = 0.9
    negative_margin: float = 0.6


class ContextualMetricLoss(nn.Module):
    def __init__(self, cfg: ContextualLossConfig):
        super().__init__()
        self.cfg = cfg

    @staticmethod
    def _mean_active(x: torch.Tensor) -> torch.Tensor:
        active = x > 0
        if not torch.any(active):
            return x.new_zeros(())
        return x[active].mean()

    def forward(self, embeddings: torch.Tensor, labels: torch.Tensor):
        W, aux = contextual_similarity(
            embeddings,
            k=self.cfg.k,
            eps=self.cfg.eps,
            alpha=self.cfg.alpha,
        )
        S = aux["S"]
        Y = labels[:, None].eq(labels[None, :]).to(S.dtype)
        off_diag = 1.0 - torch.eye(S.shape[0], device=S.device, dtype=S.dtype)

        L_context = (((W - Y) ** 2) * off_diag).mean()

        pos_hinge = F.relu(self.cfg.positive_margin - S) * Y * off_diag
        neg_hinge = F.relu(S - self.cfg.negative_margin) * (1.0 - Y)
        L_contrast = self._mean_active(pos_hinge) + self._mean_active(neg_hinge)

        L_reg = (S.mean() - self.cfg.target_mean_similarity).square()
        total = (
            self.cfg.lam * L_context
            + (1.0 - self.cfg.lam) * L_contrast
            + self.cfg.gamma * L_reg
        )
        stats = {
            "loss": total.detach(),
            "context": L_context.detach(),
            "contrast": L_contrast.detach(),
            "reg": L_reg.detach(),
            "mean_similarity": S.mean().detach(),
        }
        return total, stats
\end{lstlisting}

\section{Training loop for pose sequences}

\begin{lstlisting}
for pose, valid_mask, label in train_loader:
    pose = pose.to(device)          # [n, T, J, C]
    valid_mask = valid_mask.to(device)
    label = label.to(device)

    h = motion_encoder(pose, valid_mask=valid_mask)
    z = projection_head(h)

    loss, stats = metric_loss(z, label)

    optimizer.zero_grad(set_to_none=True)
    loss.backward()
    torch.nn.utils.clip_grad_norm_(parameters, max_norm=5.0)
    optimizer.step()
\end{lstlisting}

\section{Sampler invariant}

Every training batch should satisfy
\begin{equation}
  \#\{i:y_i=c\}=k
  \quad\text{for every class }c\text{ represented in the batch}.
  \label{eq:sampler-invariant}
\end{equation}
Use assertions during development:

\begin{lstlisting}
unique, counts = torch.unique(labels, return_counts=True)
assert torch.all(counts == k), (unique, counts)
assert batch_size % k == 0
assert k % 2 == 0
\end{lstlisting}

\section{Complexity and memory}

The dominant contextual operations multiply $n\times n$ matrices, so arithmetic scales as $O(n^3)$ with batch size. The paper also describes empirical time and space scaling as cubic, while emphasizing that GPU matrix multiplication makes the observed cost practical at its studied batch sizes. For pose models, the encoder may dominate compute, but long sequences can make batch size the binding constraint. Profile rather than assume.

Practical levers:
\begin{itemize}[leftmargin=*,itemsep=2pt]
  \item start with $16\times4=64$ or $32\times4=128$ batches;
  \item shorten or subsample sequences before shrinking $k$ below four;
  \item use mixed precision for the encoder, but inspect neighborhood ties and numerical stability;
  \item gradient accumulation does \emph{not} enlarge the contextual neighborhood, because each micro-batch forms a separate graph;
  \item parameterize $k$ explicitly---the public 256-pixel reference path hard-codes $k=4$.
\end{itemize}

\begin{warningbox}[title=Important implementation trap]
Gradient accumulation increases effective optimizer batch size but does not reproduce a larger contextual batch. The loss only sees the current micro-batch's $n\times n$ graph. If neighborhood diversity matters, you need an actual larger forward batch, a memory-bank extension, or a changed method.
\end{warningbox}

\section{Unit tests before training}

\begin{enumerate}[leftmargin=*,itemsep=4pt]
  \item \textbf{Symmetry:} verify $W\approx W^\transpose$.
  \item \textbf{Range:} verify $0\le W_{ij}\le1$ up to floating error.
  \item \textbf{Perfect blocks:} construct embeddings with well-separated class clusters and check near-zero contextual loss.
  \item \textbf{Broken pair:} move one negative into an anchor's top-$k$ set and check the loss increases.
  \item \textbf{Gradient existence:} call backward and verify finite, nonzero gradients.
  \item \textbf{Sampler:} assert exact $k$ counts per represented class.
  \item \textbf{Deterministic corruptions:} fixed seed and sample ID must reproduce identical corrupted queries.
\end{enumerate}

% ================================================================
\chapter{Retrieval Evaluation and a Fair Experimental Design}
\label{chap:evaluation}
% ================================================================

\section{Begin with one ranked list}

Suppose a query has three relevant gallery items. The system returns the following top six labels:
\[
  [\checkmark,\times,\checkmark,\times,\times,\checkmark].
\]
Recall@1 asks whether the first result is relevant. Recall@5 asks whether at least one relevant item appears among the first five. Average precision rewards a ranking that places all relevant items early, not merely one relevant item.

\begin{conceptbox}[title=Training loss and evaluation metric are different]
The model is optimized using a differentiable training objective. It is judged using retrieval metrics computed from complete ranked lists. A useful training loss is a surrogate: minimizing it should improve the non-differentiable evaluation metrics.
\end{conceptbox}



\section{Query/gallery construction}

Use disjoint query, gallery, and validation partitions. The query clip must not occur in the gallery. The primary deployment model is
\[
\text{corrupted query}\quad\longrightarrow\quad\text{fixed clean gallery}.
\]
A secondary experiment may corrupt both sides, but it answers a different question.

\section{Recall at $K$}

For query $q$, let $\pi_q(r)$ be the gallery item at rank $r$. Then
\begin{equation}
  \operatorname{R@K}(q)=
  \ind\left[\exists r\le K:\operatorname{rel}(q,\pi_q(r))=1\right].
  \label{eq:recallk}
\end{equation}
Dataset R@K is the mean over queries.

\section{Average precision and mAP}

Let $R_q$ be the number of relevant gallery items for $q$, and let $\operatorname{Prec}_q(r)$ be precision among the first $r$ retrieved items. Then
\begin{equation}
  \operatorname{AP}(q)=\frac{1}{R_q}
  \sum_{r=1}^{m}\operatorname{Prec}_q(r)\,
  \operatorname{rel}(q,\pi_q(r)),
  \label{eq:ap}
\end{equation}
with
\begin{equation}
  \operatorname{mAP}=\frac{1}{\card{\mathcal Q}}
  \sum_{q\in\mathcal Q}\operatorname{AP}(q).
  \label{eq:map}
\end{equation}
R@1 asks whether the first answer is useful; mAP evaluates the ordering of all relevant examples.

\section{Robustness curves}

Report the raw metric and the clean-to-corrupted drop. A useful summary beyond the proposal is the normalized retention
\begin{equation}
  \rho_\ell(c)=\frac{M_\ell(c)}{M_\ell(0)+\eta}.
  \label{eq:retention}
\end{equation}
This separates clean strength from robustness. It should supplement, not replace, absolute metrics.

\begin{figure}[H]
\centering
\begin{tikzpicture}
\begin{axis}[
  width=0.78\textwidth,height=6.2cm,xlabel={corruption severity},ylabel={retrieval metric},
  xmin=0,xmax=4,ymin=0.3,ymax=0.95,xtick={0,1,2,3,4},xticklabels={clean,mild,medium,strong,severe},
  grid=both,legend style={at={(0.5,-0.22)},anchor=north,legend columns=3}]
  \addplot[navy,very thick,mark=*] coordinates {(0,0.86)(1,0.82)(2,0.76)(3,0.67)(4,0.55)};
  \addplot[teal,very thick,mark=square*] coordinates {(0,0.84)(1,0.82)(2,0.79)(3,0.74)(4,0.66)};
  \addplot[gold,very thick,mark=triangle*] coordinates {(0,0.87)(1,0.84)(2,0.80)(3,0.73)(4,0.62)};
  \legend{SupCon,contextual,hybrid}
\end{axis}
\end{tikzpicture}
\caption{Schematic robustness curves only---not experimental results. The target result is a smaller degradation slope, not necessarily the highest clean point.}
\label{fig:robustness-schematic}
\end{figure}

\section{Core comparison matrix}

\begin{center}
\begin{tabularx}{\textwidth}{>{\bfseries}l X X X}
\toprule
Method & Same encoder? & Same corruption exposure? & Scientific role \\
\midrule
cross-entropy embedding & yes & matched & classification baseline \\
triplet or multi-similarity & yes & matched & conventional metric learning \\
supervised contrastive & yes & matched & strong batch-contrastive baseline \\
contextual only & yes & matched & isolate neighborhood objective \\
hybrid contextual+contrastive & yes & matched & paper's complete mechanism \\
corruption-aware SupCon & yes & same augmented inputs & separate loss effect from augmentation \\
MaskCLR-style control & preferably yes & protocol-matched & robustness-specific comparator \\
\bottomrule
\end{tabularx}
\end{center}

\section{Fairness controls}

Following the metric-learning methodology cautions in Musgrave et al., hold fixed:
\begin{itemize}[leftmargin=*,itemsep=2pt]
  \item encoder architecture and initialization;
  \item embedding dimension and normalization;
  \item data split and query/gallery manifest;
  \item augmentation and corruption exposure;
  \item number of updates and tuning budget;
  \item batch composition, except where a method inherently requires a different sampler;
  \item model-selection metric and validation set;
  \item reporting over multiple seeds.
\end{itemize}

\section{Recommended experiment order}

\begin{enumerate}[leftmargin=*,itemsep=4pt]
  \item clean frozen-feature retrieval;
  \item corruption curves with no training;
  \item projection-head SupCon vs contextual vs hybrid;
  \item full-encoder fine-tuning;
  \item corruption-aware training control;
  \item cross-subject and cross-setup replication;
  \item optional pose-estimator shift;
  \item optional SportsPose or fine-grained sports extension.
\end{enumerate}

% ================================================================
\chapter{Discussion Preparation}
\label{chap:meeting}
% ================================================================

\section{A 60-second explanation}

\begin{meetingbox}[title=Use this almost verbatim]
``I want to test whether contextual metric learning transfers from image retrieval to pose-sequence retrieval. A temporal encoder maps each skeleton sequence to a normalized vector, and cosine similarity ranks a clean gallery. Standard metric losses supervise individual pairs or triplets. The contextual objective also compares their local neighborhoods: it builds a differentiable top-$k$ graph, measures shared and reciprocal neighbors, and trains that graph to match label agreement. The original paper shows robustness to label noise, but that does not imply robustness to corrupted pose coordinates. My experiment isolates that question by holding the encoder, data split, batch composition, and corruption protocol fixed while comparing supervised contrastive, contextual, and hybrid losses on clean and increasingly corrupted queries.''
\end{meetingbox}

\section{The five equations to narrate}

\begin{enumerate}[leftmargin=*,itemsep=7pt]
  \item $D(i,j)=2-2s_{ij}$: normalized cosine geometry.
  \item $N_{ij}=\theta(-D_{ij}+\sg(D_{i,k})+\epsilon)$: differentiable neighborhood membership.
  \item $W_1$: normalized agreement on shared neighbors and shared non-neighbors.
  \item $W_2$ and $W$: reciprocal query expansion and symmetrization.
  \item $\mathcal L=\lambda\mathcal L_{context}+(1-\lambda)\mathcal L_{contrast}+\gamma\mathcal L_{reg}$: ranking, margins, and space usage.
\end{enumerate}

\section{Likely questions and precise answers}

\begin{longtable}{p{0.31\textwidth}p{0.62\textwidth}}
\toprule
\textbf{Question} & \textbf{Answer} \\
\midrule
What is the novelty? & Not pose retrieval itself and not a new encoder. The novelty is a controlled test of whether a supervised neighborhood-ranking objective improves \emph{retrieval robustness to corrupted pose inputs}. \\
Why should context help input noise? & It may reduce reliance on one unstable pairwise comparison by aggregating local neighborhood evidence. That is a hypothesis; coherent neighborhood shifts could also make it fail. \\
Why NTU120 first? & It provides scale, many action classes, standard subject/setup protocols, and compatibility with a public motion-embedding pipeline. It is an engineering and methodology benchmark, not proof of fine-grained technique understanding. \\
Why $k$ equals samples per class? & Under balanced sampling, perfect top-$k$ neighborhoods then equal the same-label set. This is the condition behind the ranking proposition. \\
Why not only use SupCon? & SupCon is the strongest simple baseline. The project asks whether explicit neighborhood overlap adds anything beyond many-positive contrastive learning. \\
What is the cleanest primary endpoint? & The difference in degradation curves under query corruption, with the gallery fixed and clean, reported using R@1/5/10 and mAP. \\
What would a null result mean? & Robustness comes mainly from encoder pretraining or corruption exposure rather than the neighborhood objective. That is still a useful separation of mechanisms. \\
What is the main technical risk? & Batch-level cubic work and the need for enough samples per class. Profile $16\times4$ and $32\times4$ batches before full experiments. \\
Could action-class relevance be too coarse? & Yes. It is appropriate for the first falsifiable benchmark, but it cannot support claims about movement quality, technique, or rehabilitation status. \\
How will you ensure fairness? & Same backbone, initialization, split, sampler, augmentation, tuning budget, and evaluator; only the objective changes in the core comparison. \\
\bottomrule
\end{longtable}

\section{Three decisions to ask the professor}

\begin{enumerate}[leftmargin=*,itemsep=5pt]
  \item Is NTU class-level retrieval an appropriate first scientific target, or should relevance be fine-grained from the outset?
  \item Should the first deliverable be a single-backbone loss comparison, or include a full MaskCLR-style robustness baseline immediately?
  \item Is the strongest framing ``contextual metric learning for temporal data'' or ``retrieval robustness under pose-estimation noise''?
\end{enumerate}

\section{Claims to avoid}

\begin{warningbox}[title=Do not say these]
\begin{itemize}[leftmargin=*,itemsep=2pt]
  \item ``Contextual loss is robust to pose noise.'' It has not yet been tested.
  \item ``NTU retrieval measures technique quality.'' It measures class relevance under the chosen protocol.
  \item ``The project invents a new motion architecture.'' The initial strength is controlled integration, not architecture novelty.
  \item ``A higher clean R@1 proves robustness.'' Robustness requires degradation curves or cross-domain evaluation.
  \item ``Gradient accumulation solves the batch-size problem.'' It does not create a larger contextual graph.
\end{itemize}
\end{warningbox}

\section{Final mental model}

\begin{conceptbox}[title=One sentence per layer]
\begin{enumerate}[leftmargin=*,itemsep=3pt]
  \item \textbf{Geometry:} normalized vectors turn cosine ranking into distance ranking.
  \item \textbf{Graph:} top-$k$ comparisons turn the similarity matrix into a local directed graph.
  \item \textbf{Context:} shared neighbors and non-neighbors compare graph signatures.
  \item \textbf{Consensus:} reciprocal query expansion smooths scores through trusted neighbors.
  \item \textbf{Learning:} MSE matches the contextual graph to label agreement, while contrastive margins and a mean-similarity regularizer stabilize the embedding.
  \item \textbf{Research question:} does that graph-based inductive bias preserve motion retrieval when the observed skeleton is damaged?
\end{enumerate}
\end{conceptbox}

% ================================================================
\appendix
\chapter{Notation Reference}
% ================================================================

\begin{longtable}{>{\bfseries}p{0.14\textwidth}p{0.75\textwidth}}
\toprule
Symbol & Meaning \\
\midrule
$x_i$ & input item; an image in the foundation paper or pose sequence in the proposal \\
$h_i$ & unnormalized encoder output \\
$z_i\in\R^d$ & unit-normalized embedding \\
$s_{ij}$ & cosine similarity $z_i^\transpose z_j$ \\
$D(i,j)$ & squared Euclidean distance $2-2s_{ij}$ \\
$n$ & mini-batch size \\
$k$ & samples per class and contextual neighborhood size \\
$y_{ij}$ & binary label similarity \\
$\theta$ & Heaviside comparison with custom backward gradient \\
$\alpha$ & constant surrogate derivative for $\theta$ \\
$\epsilon$ & neighborhood margin \\
$N_{k+\epsilon}(i)$ & expanded $k$-nearest-neighbor set of sample $i$ \\
$M_+$ & count of shared neighbors \\
$M_-$ & count of shared non-neighbors \\
$W_1$ & initial gated contextual similarity \\
$R_{k/2+\epsilon}(i)$ & reciprocal close-neighbor set \\
$W_2$ & query-expanded contextual similarity \\
$w_{ij}$ & final symmetric contextual similarity \\
$\delta_+,\delta_-$ & positive and negative contrastive margins \\
$\tilde s$ & target mean cosine similarity \\
$\lambda,\gamma$ & loss-mixture weights \\
\bottomrule
\end{longtable}

% ================================================================
\chapter{Detailed Algebra and Gradient Notes}
% ================================================================

\section{Complement-intersection simplification}

Starting from
\[
M_-(i,j)=\sum_p(1-N_{ip})(1-N_{jp}),
\]
expand:
\begin{align}
M_-(i,j)
&=\sum_p\left(1-N_{ip}-N_{jp}+N_{ip}N_{jp}\right)\\
&=n-\sum_pN_{ip}-\sum_pN_{jp}+\sum_pN_{ip}N_{jp}\\
&=n-2k+\inner{N_i}{N_j}.
\end{align}
Substituting into $W_1$ gives the constants $a$ and $b$ in \cref{eq:ab}.

\section{Step-1 surrogate gradient}

For the simplified loss
\[
\mathcal L_1=\operatorname{MSE}(y_{ij},N_{ij}),
\]
a negative sample inside $N(i)$ receives a constant push outward, while a positive sample outside $N(i)$ receives a constant pull inward. This resembles contrastive loss with a dynamic boundary: the current $k$-th-neighbor distance replaces a fixed margin.

\section{Step-2 context gradient}

For a wrongly included negative pair and $\epsilon=0$, the supplementary derivation shows that optimizing neighborhood overlap changes the distance from $i$ to every sample $p$ depending on whether $p$ lies inside $j$'s neighborhood. In compact qualitative form,
\[
\operatorname{sign}\!\left(\frac{\partial\mathcal L}{\partial D(i,p)}\right)
=
\begin{cases}
-1,&p\in N(j),\\
+1,&p\notin N(j),
\end{cases}
\]
under the paper's sign convention for the considered negative error. Gradient descent therefore increases distance to $j$'s context and decreases distance to points outside that context. The exact coefficients depend on $k$, $n-k$, $\alpha$, and the upstream MSE gradient.

\section{Logical AND as multiplication}

The reciprocal mask uses
\[
R=N_2\had N_2^\transpose.
\]
For binary inputs, multiplication equals logical AND. On the continuous backward surrogate, multiplication is smooth but has zero derivative when the other factor is zero. The paper's ablation favors multiplication over a minimum-based alternative.

% ================================================================
\chapter{Implementation Checklist}
% ================================================================

\begin{meetingbox}[title=Before the first real run]
\begin{enumerate}[leftmargin=*,itemsep=3pt]
  \item Freeze and version the query/gallery split manifest.
  \item Verify no query ID appears in the gallery.
  \item Save one deterministic corrupted copy per severity for debugging.
  \item Confirm every batch has exactly $k$ examples per represented class.
  \item Log $\mathcal L_{context}$, $\mathcal L_{contrast}$, $\mathcal L_{reg}$, mean similarity, and gradient norm separately.
  \item Visualize $S$, $N$, $W_1$, and $W$ as heatmaps for a small batch.
  \item Check clean retrieval before any corruption experiment.
  \item Compare exact brute-force cosine retrieval against any FAISS index before using approximate search.
  \item Tune each baseline on validation data with a matched budget; never select on test corruption curves.
  \item Run at least three seeds once the pipeline is stable.
\end{enumerate}
\end{meetingbox}

% ================================================================
\chapter{Two-Minute Detachable Cheat Sheet}
% ================================================================

\begin{tcolorbox}[enhanced,breakable,colback=mist,colframe=navy,title=Core proposition,fonttitle=\bfseries]
\textbf{Question:} Does contextual neighborhood optimization preserve pose-sequence retrieval better than standard metric losses as skeleton observations become corrupted?

\textbf{Design:} same MotionBERT encoder, NTU120 split, batch sampler, augmentations, and evaluator; compare SupCon, contextual, and hybrid; corrupt query only against a fixed clean gallery.
\end{tcolorbox}

\begin{tcolorbox}[enhanced,breakable,colback=sky,colframe=teal,title=Five equations,fonttitle=\bfseries]
\[
D=2-2S
\]
\[
N_{ij}=\theta(-D_{ij}+\sg(D_{i,k})+\epsilon)
\]
\[
W_1=\frac12\left(\frac{NN^\transpose}{\card{N_i}}+
\frac{(1-N)(1-N)^\transpose}{\card{N_i^c}}\right)\had N
\]
\[
W=\tfrac12\left(\frac{RW_1}{\card{R_i}}+
\left(\frac{RW_1}{\card{R_i}}\right)^\transpose\right)
\]
\[
\mathcal L=\lambda\operatorname{MSE}(W,Y)+(1-\lambda)\mathcal L_{contrast}+\gamma\mathcal L_{reg}
\]
\end{tcolorbox}

\begin{tcolorbox}[enhanced,breakable,colback=amber,colframe=gold,title=What makes the method non-standard,fonttitle=\bfseries]
Exact binary top-$k$ comparisons in the forward pass; constant heuristic gradients in the backward pass; neighborhood overlap and reciprocal query expansion; $k$ tied to samples per class.
\end{tcolorbox}

\begin{tcolorbox}[enhanced,breakable,colback=rose,colframe=brick,title=What remains unknown,fonttitle=\bfseries]
The paper's label-noise robustness does not prove pose-input robustness. The project must measure corruption degradation curves and include corruption-aware controls.
\end{tcolorbox}

% ================================================================
\backmatter
\chapter*{References}
\addcontentsline{toc}{chapter}{References}
\begin{thebibliography}{99}

\bibitem{liao2023}
C. Liao, T. Tsiligkaridis, and B. Kulis.
\newblock Supervised Metric Learning to Rank for Retrieval via Contextual Similarity Optimization.
\newblock In \emph{Proceedings of the 40th International Conference on Machine Learning}, 2023.

\bibitem{khosla2020}
P. Khosla et al.
\newblock Supervised Contrastive Learning.
\newblock In \emph{NeurIPS}, 2020.

\bibitem{musgrave2020}
K. Musgrave, S. Belongie, and S.-N. Lim.
\newblock A Metric Learning Reality Check.
\newblock In \emph{ECCV}, 2020.

\bibitem{kim2022}
S. Kim, D. Kim, M. Cho, and S. Kwak.
\newblock Self-Taught Metric Learning without Labels.
\newblock In \emph{CVPR}, 2022.

\bibitem{ntu120}
J. Liu et al.
\newblock NTU RGB+D 120: A Large-Scale Benchmark for 3D Human Activity Understanding.
\newblock \emph{IEEE TPAMI}, 2020.

\bibitem{skeletonDML}
R. Memmesheimer, S. Haring, N. Theisen, and D. Paulus.
\newblock Skeleton-DML: Deep Metric Learning for Skeleton-Based One-Shot Action Recognition.
\newblock In \emph{WACV}, 2022.

\bibitem{motionbert}
W. Zhu et al.
\newblock MotionBERT: A Unified Perspective on Learning Human Motion Representations.
\newblock In \emph{ICCV}, 2023.

\bibitem{maskclr}
M. Abdelfattah, M. Hassan, and A. Alahi.
\newblock MaskCLR: Attention-Guided Contrastive Learning for Robust Action Representation Learning.
\newblock In \emph{CVPR}, 2024.

\bibitem{faiss}
M. Douze et al.
\newblock The FAISS Library.
\newblock arXiv:2401.08281.

\bibitem{repo}
C. Liao.
\newblock Reference implementation: \emph{metric-learning-using-contextual-similarity}.
\newblock GitHub repository associated with \cite{liao2023}.

\end{thebibliography}

\end{document}
